\documentclass[11pt,twoside,openright]{book}

% Estilos, paquetes, geometría, colores y entornos pedagógicos
% ==============================================================================
% ESTILOS Y PAQUETES: LIBRO UNIVERSITARIO DE ESTADÍSTICA II (MAT-260) - UAGRM
% MSc. Ing. Anselmo Salguero Arano
% ==============================================================================

\usepackage[utf8]{inputenc}
\usepackage[T1]{fontenc}
\usepackage[spanish,es-tabla,es-nodecimaldot]{babel}
\usepackage{amsmath,amssymb,amsfonts,amsthm,mathtools}
\usepackage{lmodern}
\usepackage{geometry}
\geometry{
    letterpaper,
    top=2.5cm,
    bottom=2.8cm,
    left=3.0cm,
    right=2.5cm,
    headheight=25pt,
    footskip=1.4cm
}

% --- Microtype seguro ---
\usepackage[protrusion=true,expansion=false]{microtype}

% --- Colores Institucionales y de Diseño ---
\usepackage{xcolor}
\definecolor{UAGRMBlue}{RGB}{11, 47, 100}       % Azul institucional UAGRM
\definecolor{UAGRMRed}{RGB}{140, 20, 20}        % Rojo institucional
\definecolor{UAGRMDark}{RGB}{20, 30, 45}        % Texto principal / títulos oscuros
\definecolor{AccentGreen}{RGB}{20, 110, 60}     % Verde contable / éxito
\definecolor{AccentGold}{RGB}{180, 130, 20}     % Dorado / énfasis
\definecolor{LightBlueBg}{RGB}{242, 246, 252}   % Fondo suave azul
\definecolor{LightGreenBg}{RGB}{242, 250, 244}  % Fondo suave verde
\definecolor{LightRedBg}{RGB}{253, 244, 244}    % Fondo suave rojo/alerta
\definecolor{LightGrayBg}{RGB}{248, 249, 250}   % Fondo neutro
\definecolor{BorderGray}{RGB}{210, 215, 225}    % Borde neutro

% --- Enlaces y referencias ---
\usepackage{hyperref}
\hypersetup{
    colorlinks=true,
    linkcolor=UAGRMBlue,
    citecolor=UAGRMRed,
    urlcolor=UAGRMBlue,
    pdfauthor={MSc. Ing. Anselmo Salguero Arano},
    pdftitle={Estadística II: Probabilidades e Inferencia Estadística - UAGRM},
    pdfsubject={Texto Guía Universitario MAT-260}
}

% --- Encabezados y Pies de Página (Estilo MAT-100 UAGRM) ---
\usepackage{fancyhdr}
\pagestyle{fancy}
\fancyhf{}
\fancyhead[LE,RO]{\textcolor{UAGRMDark}{\bfseries\small\thepage}}
\fancyhead[RE]{\textcolor{UAGRMBlue}{\small\nouppercase{\leftmark}}}
\fancyhead[LO]{\textcolor{UAGRMBlue}{\small\nouppercase{\rightmark}}}
\fancyfoot[LE,RO]{\textcolor{UAGRMDark!70}{\scriptsize\textit{MSc. Ing. Anselmo Salguero Arano --- UAGRM}}}
\fancyfoot[RE,LO]{\textcolor{UAGRMDark!70}{\scriptsize\textit{Estadística II (MAT-260) --- Contaduría Pública}}}
\renewcommand{\headrulewidth}{0.6pt}
\renewcommand{\footrulewidth}{0.4pt}

\fancypagestyle{plain}{
    \fancyhf{}
    \fancyfoot[C]{\textcolor{UAGRMDark}{\small\thepage}}
    \fancyfoot[R]{\textcolor{UAGRMDark!70}{\scriptsize\textit{Estadística II --- UAGRM}}}
    \renewcommand{\headrulewidth}{0pt}
    \renewcommand{\footrulewidth}{0.4pt}
}

% --- Cajas y Entornos Pedagógicos con TColorBox ---
\usepackage[most]{tcolorbox}
\tcbuselibrary{skins,breakable}

% 1. Caja de Definición Teórica
\newtcolorbox[auto counter,number within=chapter]{definicion}[2][]{
    enhanced,
    breakable,
    colback=LightBlueBg,
    colframe=UAGRMBlue,
    coltitle=white,
    fonttitle=\bfseries\sffamily,
    title=Definición~\thetcbcounter: #2,
    arc=2.5mm,
    boxrule=1pt,
    left=4mm, right=4mm, top=3mm, bottom=3mm,
    #1
}

% 2. Caja de Ejemplo Resuelto Paso a Paso (Anderson & Lind)
\newtcolorbox[auto counter,number within=chapter]{ejemplo}[2][]{
    enhanced,
    breakable,
    colback=white,
    colframe=UAGRMBlue!85!black,
    coltitle=white,
    fonttitle=\bfseries\sffamily,
    title=Ejercicio Resuelto~\thetcbcounter: #2,
    arc=2.5mm,
    boxrule=1.2pt,
    titlerule=0.5pt,
    colbacktitle=UAGRMBlue,
    left=4mm, right=4mm, top=3.5mm, bottom=3.5mm,
    #1
}

% 3. Caja de Aplicación en Contaduría / Auditoría / Finanzas
\newtcolorbox[auto counter,number within=chapter]{casocontable}[2][]{
    enhanced,
    breakable,
    colback=LightGreenBg,
    colframe=AccentGreen,
    coltitle=white,
    fonttitle=\bfseries\sffamily,
    title=Caso Práctico en Contaduría y Finanzas~\thetcbcounter: #2,
    arc=2.5mm,
    boxrule=1pt,
    left=4mm, right=4mm, top=3mm, bottom=3mm,
    #1
}

% 4. Caja de Alerta / Error Común
\newtcolorbox{alerta}[1][¡Atención --- Error Común a Evitar!]{
    enhanced,
    breakable,
    colback=LightRedBg,
    colframe=UAGRMRed,
    coltitle=white,
    fonttitle=\bfseries\sffamily,
    title=#1,
    arc=2mm,
    boxrule=1pt,
    left=3.5mm, right=3.5mm, top=2.5mm, bottom=2.5mm
}

% 5. Caja de Formulario Resumen
\newtcolorbox{formulario}[1][Formulario Resumen de la Unidad]{
    enhanced,
    breakable,
    colback=LightGrayBg,
    colframe=UAGRMDark,
    coltitle=white,
    fonttitle=\bfseries\sffamily,
    title=#1,
    arc=2.5mm,
    boxrule=1.2pt,
    left=4mm, right=4mm, top=3mm, bottom=3mm
}

% --- Tablas y Gráficos ---
\usepackage{booktabs}
\usepackage{tabularx}
\usepackage{multirow}
\usepackage{array}
\usepackage{graphicx}
\usepackage{enumitem}
\usepackage{tikz}
\usetikzlibrary{shapes,arrows,positioning,calc,patterns}

% --- Formato de Títulos de Capítulos con Titlesec ---
\usepackage{titlesec}
\titleformat{\chapter}[display]
  {\normalfont\sffamily\huge\bfseries\color{UAGRMBlue}}
  {\flushright\normalsize\bfseries\color{UAGRMRed}UNIDAD \Huge\thechapter}
  {1ex}
  {\vspace{1ex}}

\titleformat{\section}
  {\normalfont\sffamily\Large\bfseries\color{UAGRMBlue}}
  {\thesection}{1em}{}

\titleformat{\subsection}
  {\normalfont\sffamily\large\bfseries\color{UAGRMDark}}
  {\thesubsection}{1em}{}


\begin{document}

% --- PARTE PRELIMINAR (FRONTMATTER) ---
\frontmatter

% Portada Oficial
% ==============================================================================
% PORTADA INSTITUCIONAL PROFESIONAL - UAGRM
% ==============================================================================
\begin{titlepage}
\begin{tikzpicture}[remember picture,overlay]
    % Franja lateral decorativa institucional
    \fill[color=UAGRMBlue] (current page.north west) rectangle ([xshift=1.8cm]current page.south west);
    \fill[color=UAGRMRed] ([xshift=1.8cm]current page.north west) rectangle ([xshift=2.2cm]current page.south west);
    \fill[color=AccentGold] ([xshift=2.2cm]current page.north west) rectangle ([xshift=2.35cm]current page.south west);
\end{tikzpicture}

\vspace{-1.5cm}
\begin{flushright}
    {\scshape\large\bfseries UNIVERSIDAD AUTÓNOMA GABRIEL RENÉ MORENO}\\[0.2cm]
    {\scshape\normalsize FACULTAD DE CIENCIAS CONTABLES, AUDITORÍA,}\\
    {\scshape\normalsize SISTEMAS DE CONTROL DE GESTIÓN Y FINANZAS}\\[0.2cm]
    {\scshape\small\bfseries CARRERA DE CONTADURÍA PÚBLICA}\\[0.8cm]
    \rule{0.75\textwidth}{1.5pt}\\[0.4cm]
    
    \vspace{2.5cm}
    
    {\Huge\bfseries\color{UAGRMBlue} ESTADÍSTICA II}\\[0.5cm]
    {\LARGE\bfseries\color{UAGRMRed} Probabilidades e Inferencia Estadística}\\[0.4cm]
    {\large\itshape Texto Guía Teórico-Práctico con Aplicaciones a la Contaduría, Auditoría y Finanzas}\\[0.6cm]
    {\small\sffamily\color{gray} Incluye ejercicios resueltos y casos basados en Anderson \& Lind}\\[2.0cm]
    
    \vfill
    
    \begin{minipage}{0.65\textwidth}
        \flushright
        \textbf{\large Docente Autor:}\\[0.15cm]
        {\Large\bfseries\color{UAGRMDark} MSc. Ing. Anselmo Salguero Arano}\\[0.2cm]
        {\small Catedrático de Estadística y Métodos Cuantitativos\\
        UAGRM — Santa Cruz de la Sierra, Bolivia}\\[0.8cm]
        {\bfseries\color{UAGRMBlue} Sigla de Asignatura: MAT-260}\\[0.1cm]
        {\small Edición Universitaria Revisada y Ampliada (2026)}
    \end{minipage}
\end{flushright}
\end{titlepage}

\cleardoublepage


% Página de Créditos y Presentación Institucional
\thispagestyle{empty}
\vspace*{12cm}
\begin{flushleft}
    \textbf{ESTADÍSTICA II: Probabilidades, Inferencia y Modelos de Decisión}\\
    \textit{Texto Guía Universitario Teórico-Práctico (Edición 2026)}\\[0.3cm]
    \copyright\ 2026 MSc. Ing. Anselmo Salguero Arano\\
    Universidad Autónoma Gabriel René Moreno (UAGRM)\\
    Facultad de Ciencias Contables, Auditoría, Sistemas de Control de Gestión y Finanzas\\
    Carrera de Contaduría Pública\\
    Santa Cruz de la Sierra -- Bolivia\\[0.4cm]
    \small Todos los derechos reservados. Este texto ha sido elaborado como material oficial de docencia universitaria y apoyo al aprendizaje presencial y virtual en la plataforma Moodle de la UAGRM.
\end{flushleft}
\clearpage

% Prólogo / Presentación del Autor
\chapter*{Presentación y Guía Didáctica}
\addcontentsline{toc}{chapter}{Presentación y Guía Didáctica}

El presente texto guía ha sido diseñado y estructurado específicamente para los estudiantes de la Carrera de \textbf{Contaduría Pública} de la Universidad Autónoma Gabriel René Moreno (UAGRM), respondiendo fielmente a las exigencias del programa analítico oficial de la asignatura \textbf{Estadística II (MAT-260)} y ampliando su alcance hacia la inferencia y los modelos cuantitativos contemporáneos.

En el ejercicio profesional de la contaduría pública, la auditoría tributaria y financiera, el peritaje judicial y la gestión financiera corporativa, las decisiones estratégicas se enfrentan cotidianamente a escenarios de incertidumbre y grandes volúmenes de datos transaccionales. Las estimaciones contables, la determinación del riesgo crediticio, el muestreo de comprobantes y la segregación de costos ya no pueden basarse en la simple intuición, sino en el rigor metódico de las probabilidades, la inferencia estadística y la econometría aplicada.

\section*{Estructura y Características Distintivas de la Edición 2026:}
\begin{itemize}
    \item \textbf{Siete (7) Unidades Temáticas Integrales:}
    \begin{enumerate}
        \item \textbf{Unidad 1: Análisis Combinatorio y Métodos de Conteo:} Fundamentos de conteo, permutaciones y combinaciones aplicadas a códigos contables y muestras de auditoría.
        \item \textbf{Unidad 2: Teoría y Fundamentos de Probabilidad:} Reglas de adición, multiplicación, tablas de contingencia y Teorema de Bayes en detección de fraudes y riesgo.
        \item \textbf{Unidad 3: Distribuciones de Probabilidad Discretas:} Modelado con distribuciones Binomial, Hipergeométrica y Poisson en control de calidad y auditoría de cumplimiento.
        \item \textbf{Unidad 4: Distribuciones de Probabilidad Continuas:} Distribución Normal Gaussiana, estandarización $Z$, percentiles y tiempos de servicio con la distribución Exponencial.
        \item \textbf{Unidad 5: Distribuciones Muestrales e Intervalos de Confianza:} Teorema del Límite Central, error estándar, intervalos con $Z$ y $t$ de Student, y tamaño óptimo de muestra (NIA 530).
        \item \textbf{Unidad 6: Pruebas de Hipótesis para una Población:} Toma de decisiones basada en 5 pasos para medias y proporciones, valores críticos y enfoque del valor $p$ ($p$-value).
        \item \textbf{Unidad 7: Regresión Lineal Simple y Correlación:} Diagramas de dispersión, coeficiente de Pearson $r$, ajuste por MCO, coeficiente de determinación $R^2$ y separación de costos fijos y variables.
    \end{enumerate}
    \item \textbf{Enfoque Aplicado a la Realidad Empresarial Boliviana:} La teoría matemática se fundamenta de forma rigurosa y se complementa de inmediato con casos prácticos de auditoría impositiva, finanzas corporativas y contabilidad gerencial.
    \item \textbf{Abundancia de Ejercicios Resueltos y Propuestos Graduados:} Se articulan los mejores problemas clásicos de los textos de referencia internacional más prestigiosos (\textit{Anderson, Sweeney \& Williams} y \textit{Lind, Marchal \& Wathen}).
\end{itemize}

\vspace{1.2cm}
\begin{flushright}
    \textbf{MSc. Ing. Anselmo Salguero Arano}\\
    \textit{Docente Titular de Estadística II (MAT-260)}\\
    UAGRM -- Santa Cruz de la Sierra, Bolivia
\end{flushright}
\clearpage

% Tabla de Contenidos
\tableofcontents
\clearpage

% --- CUERPO PRINCIPAL DEL LIBRO (MAINMATTER) ---
\mainmatter

% Capítulo 1: Análisis Combinatorio
\chapter{Análisis Combinatorio y Métodos de Conteo}
\label{cap:analisis_combinatorio}

\section*{Objetivos de Aprendizaje de la Unidad}
Al finalizar el estudio de esta unidad, el estudiante de Contaduría Pública será capaz de:
\begin{enumerate}
    \item Identificar y aplicar correctamente el \textbf{Principio Fundamental del Conteo} (reglas de la multiplicación y de la adición) en problemas de decisión empresarial.
    \item Diferenciar con claridad cuándo un problema requiere el cálculo de \textbf{Permutaciones} (donde el orden de asignación es determinante) o de \textbf{Combinaciones} (donde únicamente interesa la conformación del grupo).
    \item Resolver problemas complejos de muestreo para auditoría, asignación de puestos directivos, codificación de cuentas contables y seguridad informática.
    \item Manejar con destreza el álgebra de factoriales para simplificar expresiones analíticas.
\end{enumerate}

\section{Introducción e Importancia en las Ciencias Contables y Económicas}
El análisis combinatorio es la rama de las matemáticas que estudia las diversas formas en que se pueden agrupar, ordenar y seleccionar los elementos de uno o varios conjuntos. En la administración de empresas, la auditoría y las finanzas modernas, el proceso de conteo sistemático es la piedra angular sobre la cual se construye la \textbf{Teoría de Probabilidades}.

Cuando un auditor financiero necesita determinar cuántas muestras distintas de facturas puede extraer de un lote de 5,000 comprobantes, o cuando un departamento de sistemas debe estimar la robustez de las claves de acceso al sistema contable institucional, el conteo directo uno a uno resulta impráctico e ineficiente. El análisis combinatorio provee las herramientas analíticas para cuantificar todas las posibilidades de manera exacta y rápida.

\section{Principios Fundamentales del Conteo}

\subsection{Principio de la Multiplicación (Regla del Producto)}
\begin{definicion}{Principio Fundamental de la Multiplicación}
Si un experimento, procedimiento o proceso de decisión consta de $k$ etapas sucesivas, donde la primera etapa puede realizarse de $n_1$ maneras diferentes, la segunda etapa puede realizarse de $n_2$ maneras diferentes (independientemente de cómo se haya realizado la primera), y así sucesivamente hasta la $k$-ésima etapa con $n_k$ maneras, entonces el número total de formas posibles en que puede ejecutarse todo el procedimiento conjunto es:
\begin{equation}
    N = n_1 \cdot n_2 \cdot n_3 \dotsm n_k
\end{equation}
\end{definicion}

\begin{ejemplo}{Diseño de Códigos de Cuentas Contables (Basado en Anderson)}
El departamento de contabilidad de una corporación comercial en Santa Cruz está estructurando un nuevo catálogo de cuentas contables para sus activos fijos. La codificación interna consta de tres partes consecutivas:
\begin{itemize}
    \item \textbf{Parte 1 (Categoría):} Una letra del alfabeto (de entre 5 categorías posibles: A, B, C, D, E).
    \item \textbf{Parte 2 (Sucursal):} Un dígito numérico del 1 al 4 (que identifica una de las 4 sucursales operativas).
    \item \textbf{Parte 3 (Identificador del Activo):} Dos dígitos numéricos del 0 al 9, donde se permite la repetición de dígitos.
\end{itemize}
¿Cuántos códigos de activos fijos distintos se pueden generar bajo este esquema?

\tcbline
\textbf{Solución Paso a Paso:}
\begin{enumerate}
    \item \textbf{Etapa 1 (Categoría):} Hay $n_1 = 5$ opciones de letras.
    \item \textbf{Etapa 2 (Sucursal):} Hay $n_2 = 4$ opciones de dígitos.
    \item \textbf{Etapa 3 (Primer dígito del activo):} Hay $n_3 = 10$ opciones ($0, 1, \dots, 9$).
    \item \textbf{Etapa 4 (Segundo dígito del activo):} Hay $n_4 = 10$ opciones ($0, 1, \dots, 9$).
\end{enumerate}
Aplicando el Principio de la Multiplicación:
\[
N = n_1 \cdot n_2 \cdot n_3 \cdot n_4 = 5 \cdot 4 \cdot 10 \cdot 10 = 2{,}000 \text{ códigos diferentes.}
\]
\textbf{Interpretación Contable:} El sistema permite registrar hasta 2,000 activos individuales sin riesgo de duplicidad de códigos.
\end{ejemplo}

\subsection{Principio de la Adición (Regla de la Suma)}
\begin{definicion}{Principio de la Adición para Eventos Mutuamente Excluyentes}
Si una tarea o elección puede realizarse mediante la opción $A$ de $n$ maneras diferentes, o mediante la opción alternativa $B$ de $m$ maneras diferentes, y ambas opciones son \textbf{mutuamente excluyentes} (es decir, no pueden ejecutarse simultáneamente en el mismo intento), entonces el número total de formas en que puede cumplirse la tarea es:
\begin{equation}
    N = n + m
\end{equation}
En general, para $k$ alternativas mutuamente excluyentes:
\begin{equation}
    N = n_1 + n_2 + \dots + n_k
\end{equation}
\end{definicion}

\begin{ejemplo}{Selección de Proveedor de Software de Auditoría (Basado en Lind)}
Una firma de auditoría en Bolivia desea adquirir una licencia de software para la gestión de papeles de trabajo. Puede optar por proveedores nacionales, que ofrecen 3 soluciones certificadas, o por proveedores internacionales que ofrecen 5 soluciones homologadas. Si la firma únicamente adquirirá un solo paquete de software, ¿de cuántas formas puede tomar su decisión de compra?

\tcbline
\textbf{Solución:}
Puesto que la adquisición de un software nacional excluye la compra simultánea de uno internacional en esta decisión:
\[
N = n_{\text{nacionales}} + n_{\text{internacionales}} = 3 + 5 = 8 \text{ opciones disponibles.}
\]
\end{ejemplo}

\section{Factorial de un Número y sus Propiedades}
\begin{definicion}{Factorial de un Entero Positivo}
Para todo número entero positivo $n$, el factorial de $n$, denotado por $n!$, se define como el producto de todos los números enteros consecutivos desde 1 hasta $n$:
\begin{equation}
    n! = n \cdot (n-1) \cdot (n-2) \dotsm 3 \cdot 2 \cdot 1
\end{equation}
Por convención axiomática fundamental:
\begin{equation}
    0! = 1
\end{equation}
\end{definicion}

\subsection{Propiedades Fundamentales de los Factoriales}
\begin{enumerate}
    \item \textbf{Propiedad Recursiva:} $n! = n \cdot (n-1)!$
    \item \textbf{Propiedad de Cocientes:} Para $n > r$:
    \[
    \frac{n!}{r!} = n \cdot (n-1) \cdot (n-2) \dotsm (r+1)
    \]
    \item \textbf{Valores frecuentes:}
    \begin{center}
    \begin{tabular}{ccccc}
    \toprule
    $0! = 1$ & $1! = 1$ & $2! = 2$ & $3! = 6$ & $4! = 24$ \\
    $5! = 120$ & $6! = 720$ & $7! = 5{,}040$ & $8! = 40{,}320$ & $9! = 362{,}880$ \\
    \bottomrule
    \end{tabular}
    \end{center}
\end{enumerate}

\section{Permutaciones}
Una permutación es cualquier arreglo u ordenación de un conjunto de objetos donde \textbf{el orden en que se colocan es fundamental}.

\subsection{Permutaciones Lineales Simples de $n$ Objetos}
El número de formas en que se pueden ordenar $n$ objetos distintos en una fila es:
\begin{equation}
    P_n = n!
\end{equation}

\subsection{Permutaciones de $n$ Objetos Tomados de $r$ en $r$ ($nPr$)}
\begin{definicion}{Permutaciones de $n$ en $r$}
El número de arreglos ordenados de tamaño $r$ que se pueden seleccionar a partir de un conjunto de $n$ objetos distintos ($r \le n$) es:
\begin{equation}
    {}_n P_r = P(n, r) = \frac{n!}{(n-r)!} = n \cdot (n-1) \cdot (n-2) \dotsm (n-r+1)
\end{equation}
\end{definicion}

\begin{ejemplo}{Asignación de Cargos en el Directorio Financiero (Basado en Anderson)}
En una junta general de accionistas de una empresa industrial, se postulan 8 candidatos calificados para ocupar tres cargos jerárquicos: Presidente del Directorio, Vicepresidente Ejecutivo y Síndico Financiero. ¿De cuántas formas distintas pueden asignarse estos 3 puestos si ningún candidato puede ocupar más de un cargo?

\tcbline
\textbf{Solución:}
Dado que el orden y la jerarquía de los cargos son determinantes (no es lo mismo ser Presidente que Síndico), se trata de una permutación de $n = 8$ candidatos tomados de $r = 3$ en $3$:
\[
{}_8 P_3 = \frac{8!}{(8-3)!} = \frac{8!}{5!} = \frac{8 \cdot 7 \cdot 6 \cdot 5!}{5!} = 8 \cdot 7 \cdot 6 = 336 \text{ formas distintas.}
\]
\end{ejemplo}

\subsection{Permutaciones con Repetición (Elementos Indistinguibles)}
\begin{definicion}{Permutaciones con Elementos Repetidos}
Si se dispone de $n$ objetos, de los cuales $r_1$ son de un primer tipo (idénticos entre sí), $r_2$ son de un segundo tipo, \dots, y $r_k$ son de un $k$-ésimo tipo, tales que $r_1 + r_2 + \dots + r_k = n$, el número de permutaciones distinguibles de los $n$ objetos es:
\begin{equation}
    P_n^{r_1, r_2, \dots, r_k} = \frac{n!}{r_1! \cdot r_2! \dotsm r_k!}
\end{equation}
\end{definicion}

\begin{ejemplo}{Codificación de Lotes de Producción (Basado en Lind)}
El auditor de costos de una fábrica textil en el Parque Industrial de Santa Cruz revisa una serie de etiquetas de control de calidad. Una clave de lote está formada por la secuencia de 9 caracteres compuesta por: 4 letras \textbf{A}, 3 letras \textbf{B} y 2 letras \textbf{C}. ¿Cuántas secuencias de etiquetas diferentes pueden formarse?

\tcbline
\textbf{Solución:}
Tenemos un total de $n = 9$ elementos, con $r_1 = 4$ letras A, $r_2 = 3$ letras B y $r_3 = 2$ letras C:
\[
P_9^{4, 3, 2} = \frac{9!}{4! \cdot 3! \cdot 2!} = \frac{362{,}880}{24 \cdot 6 \cdot 2} = \frac{362{,}880}{288} = 1{,}260 \text{ secuencias distintas.}
\]
\end{ejemplo}

\subsection{Permutaciones Circulares}
Cuando los objetos se disponen en círculo o en una mesa redonda, no existe un primer ni un último lugar absoluto; las posiciones son relativas. Fijando un elemento como referencia, el número de permutaciones circulares de $n$ objetos es:
\begin{equation}
    P_c(n) = (n-1)!
\end{equation}

\begin{ejemplo}{Mesa de Negociación de Crédito Bancario}
Seis ejecutivos de finanzas de una empresa y del banco se reúnen alrededor de una mesa redonda para negociar las condiciones de un fideicomiso. ¿De cuántas maneras diferentes pueden sentarse alrededor de la mesa?

\tcbline
\textbf{Solución:}
\[
P_c(6) = (6-1)! = 5! = 120 \text{ ordenaciones posibles.}
\]
\end{ejemplo}

\section{Combinaciones}
\begin{definicion}{Combinaciones Simples ($nCr$)}
Una combinación es una selección de $r$ objetos tomados a partir de un conjunto de $n$ objetos distintos ($r \le n$), \textbf{sin importar el orden de colocación ni la secuencia de selección}. El número de combinaciones está dado por:
\begin{equation}
    {}_n C_r = \binom{n}{r} = \frac{n!}{r! \cdot (n-r)!}
\end{equation}
\end{definicion}

\subsection{Propiedades Clave de las Combinaciones}
\begin{enumerate}
    \item \textbf{Propiedad de Simetría:} $\binom{n}{r} = \binom{n}{n-r}$
    \item \textbf{Casos Extremos:} $\binom{n}{0} = 1, \quad \binom{n}{n} = 1, \quad \binom{n}{1} = n$
    \item \textbf{Relación entre Permutaciones y Combinaciones:}
    \begin{equation}
        {}_n P_r = r! \cdot {}_n C_r \iff {}_n C_r = \frac{{}_n P_r}{r!}
    \end{equation}
\end{enumerate}

\begin{alerta}
\textbf{¿Cómo saber si usar Permutación o Combinación?}
\begin{itemize}
    \item Si al cambiar el orden de los elementos seleccionados se obtiene un resultado o asignación \textbf{diferente} (ej. Presidente vs. Secretario, orden de llegada, puestos con distinta función) $\implies$ \textbf{PERMUTACIÓN}.
    \item Si al cambiar el orden de los elementos el grupo sigue siendo exactamente \textbf{el mismo} (ej. muestra de facturas a auditar, comisión de trabajo, mano de cartas) $\implies$ \textbf{COMBINACIÓN}.
\end{itemize}
\end{alerta}

\begin{ejemplo}{Muestreo de Auditoría Financiera (Basado en Anderson)}
Un equipo de auditoría externa debe inspeccionar una muestra aleatoria de 5 cuentas de clientes morosos a partir de una lista que contiene 12 cuentas críticas en mora.
\begin{enumerate}[label=\alph*)]
    \item ¿Cuántas muestras distintas de 5 cuentas pueden seleccionarse?
    \item Si la lista contiene 7 cuentas comerciales y 5 cuentas industriales, ¿cuántas muestras de 5 cuentas contendrán exactamente 3 cuentas comerciales y 2 industriales?
\end{enumerate}

\tcbline
\textbf{Solución:}
\begin{enumerate}[label=\alph*)]
    \item \textbf{Muestras totales de 5 cuentas de entre 12:}
    Como el orden en que se extraen las cuentas para su revisión no altera el conjunto auditado:
    \[
    \binom{12}{5} = \frac{12!}{5! \cdot (12-5)!} = \frac{12 \cdot 11 \cdot 10 \cdot 9 \cdot 8}{5 \cdot 4 \cdot 3 \cdot 2 \cdot 1} = \frac{95{,}040}{120} = 792 \text{ muestras posibles.}
    \]
    \item \textbf{Muestras con 3 comerciales y 2 industriales:}
    Se descompone en dos selecciones independientes aplicando la regla de la multiplicación:
    \begin{itemize}
        \item Formas de elegir 3 comerciales de entre 7: $\binom{7}{3} = \frac{7 \cdot 6 \cdot 5}{3 \cdot 2 \cdot 1} = 35$
        \item Formas de elegir 2 industriales de entre 5: $\binom{5}{2} = \frac{5 \cdot 4}{2 \cdot 1} = 10$
    \end{itemize}
    Total de muestras con la condición requerida:
    \[
    N = \binom{7}{3} \cdot \binom{5}{2} = 35 \cdot 10 = 350 \text{ muestras.}
    \]
\end{enumerate}
\end{ejemplo}

\begin{ejemplo}{Conformación de un Tribunal de Grado en Contaduría (Basado en Lind)}
El decanato de la Facultad de Ciencias Contables debe designar un tribunal calificador para una defensa de tesis de grado integrado por 4 docentes. Se dispone de 6 docentes del área de Auditoría y 5 docentes del área de Tributación. ¿De cuántas formas se puede conformar el tribunal si debe estar integrado por al menos 2 docentes del área de Auditoría?

\tcbline
\textbf{Solución:}
La condición ``al menos 2 de Auditoría'' en un tribunal de 4 miembros abarca tres casos mutuamente excluyentes:
\begin{itemize}
    \item \textbf{Caso 1 (Exactamente 2 de Auditoría y 2 de Tributación):}
    \[
    N_1 = \binom{6}{2} \cdot \binom{5}{2} = 15 \cdot 10 = 150
    \]
    \item \textbf{Caso 2 (Exactamente 3 de Auditoría y 1 de Tributación):}
    \[
    N_2 = \binom{6}{3} \cdot \binom{5}{1} = 20 \cdot 5 = 100
    \]
    \item \textbf{Caso 3 (Exactamente 4 de Auditoría y 0 de Tributación):}
    \[
    N_3 = \binom{6}{4} \cdot \binom{5}{0} = 15 \cdot 1 = 15
    \]
\end{itemize}
Aplicando el Principio de la Adición:
\[
N_{\text{total}} = N_1 + N_2 + N_3 = 150 + 100 + 15 = 265 \text{ formas distintas de integrar el tribunal.}
\]
\end{ejemplo}

\section{Formulario Resumen del Capítulo 1}
\begin{formulario}[Resumen de Fórmulas — Análisis Combinatorio]
\begin{itemize}
    \item \textbf{Regla del Producto:} $N = n_1 \cdot n_2 \dotsm n_k$
    \item \textbf{Regla de la Suma (Excluyentes):} $N = n_1 + n_2 + \dots + n_k$
    \item \textbf{Permutaciones Lineales Simples:} $P_n = n!$
    \item \textbf{Permutaciones de $n$ en $r$ ($nPr$):} ${}_n P_r = \dfrac{n!}{(n-r)!}$
    \item \textbf{Permutaciones con Repetición:} $P_n^{r_1, r_2, \dots, r_k} = \dfrac{n!}{r_1! \cdot r_2! \dotsm r_k!}$
    \item \textbf{Permutaciones Circulares:} $P_c(n) = (n-1)!$
    \item \textbf{Combinaciones Simples ($nCr$):} ${}_n C_r = \dbinom{n}{r} = \dfrac{n!}{r!(n-r)!}$
\end{itemize}
\end{formulario}

\section{Ejercicios Propuestos de la Unidad}
\subsection*{Nivel 1: Principios Fundamentales de Conteo}
\begin{enumerate}
    \item Una microempresa de consultoría contable cuenta con 9 profesionales. Se desea conformar una comisión de 3 consultores para realizar una auditoría ambiental. ¿De cuántas formas se puede seleccionar la comisión?
    \item Un sistema informático bancario requiere que cada usuario cree un PIN de 4 dígitos. ¿Cuántos PINs diferentes pueden crearse si:
    \begin{enumerate}[label=\alph*)]
        \item Se permite repetir dígitos?
        \item No se permite repetir ningún dígito?
        \item El primer dígito no puede ser 0 y no se permiten repeticiones?
    \end{enumerate}
    \item Un analista de inversiones desea estructurar un portafolio seleccionando 2 bonos soberanos de entre 5 opciones disponibles y 3 acciones corporativas de entre 8 opciones disponibles. ¿Cuántos portafolios distintos puede diseñar?
\end{enumerate}

\subsection*{Nivel 2: Permutaciones y Muestreo en Auditoría}
\begin{enumerate}[resume]
    \item En una empresa importadora, se inspeccionan 15 contenedores de mercadería, de los cuales 4 presentan observaciones aduaneras y 11 están en regla. Si un fiscalizador elige al azar una muestra de 4 contenedores, ¿de cuántas formas puede contener la muestra:
    \begin{enumerate}[label=\alph*)]
        \item Exactamente 2 contenedores con observaciones?
        \item Ningún contenedor con observaciones?
        \item Al menos 1 contenedor con observaciones?
    \end{enumerate}
    \item Se deben asignar 5 auditores a 5 sucursales bancarias diferentes en Santa Cruz (una sucursal por auditor). ¿De cuántas maneras se puede realizar la asignación si el auditor senior necesariamente debe ser asignado a la sucursal central?
    \item ¿Cuántas palabras con o sin sentido pueden formarse con todas las letras de la palabra «CONTABILIDAD»?
\end{enumerate}

\subsection*{Nivel 3: Casos Integrales de Aplicación Empresarial}
\begin{enumerate}[resume]
    \item Una junta general de accionistas de una empresa agroindustrial en Santa Cruz debe elegir una mesa directiva de 4 miembros: Presidente, Vicepresidente, Secretario y Tesorero. Se postulan 12 candidatos (7 del sector accionista mayoritario y 5 del sector minoritario).
    \begin{enumerate}[label=\alph*)]
        \item ¿Cuántas mesas directivas distintas se pueden conformar en total?
        \item ¿En cuántas de ellas el Presidente es del sector mayoritario y el Tesorero del sector minoritario?
    \end{enumerate}
    \item Un auditor fiscal revisa un lote de 20 facturas comerciales, de las cuales 6 presentan irregularidades en el cálculo del crédito fiscal IVA. Si se extrae una muestra de 5 facturas para un peritaje judicial, ¿de cuántas maneras la muestra contendrá a lo sumo 2 facturas con irregularidades?
\end{enumerate}


% Capítulo 2: Probabilidades
\chapter{Teoría y Fundamentos de Probabilidad}
\label{cap:probabilidades}

\section*{Objetivos de Aprendizaje de la Unidad}
Al finalizar el estudio de esta unidad, el estudiante será capaz de:
\begin{enumerate}
    \item Comprender los fundamentos de la incertidumbre en los negocios y diferenciar entre los enfoques \textbf{clásico}, de \textbf{frecuencia relativa} y \textbf{subjetivo} de probabilidad.
    \item Aplicar con exactitud los axiomas de probabilidad y las reglas de la \textbf{adición} y de la \textbf{multiplicación} para eventos dependientes e independientes.
    \item Construir e interpretar \textbf{tablas de contingencia bidimensionales} para calcular probabilidades conjuntas, marginales y condicionales en análisis de mercado y auditoría.
    \item Dominar y aplicar el \textbf{Teorema de Bayes} para actualizar probabilidades a priori con nueva información en análisis de riesgo crediticio, fraude contable y toma de decisiones gerenciales.
\end{enumerate}

\section{Introducción a la Probabilidad en el Entorno Contable-Empresarial}
En el mundo de los negocios, la contabilidad y las finanzas, casi ninguna decisión directiva se toma bajo condiciones de certeza absoluta. Un auditor no revisa el 100\% de las transacciones de una empresa, sino una muestra; un oficial de crédito no puede predecir con certeza si un prestatario pagará a tiempo; y un analista financiero debe evaluar proyectos con flujos de caja sujetos a riesgo de mercado.

La \textbf{Probabilidad} es el lenguaje matemático que permite medir y modelar cuantitativamente la incertidumbre, permitiendo transformar corazonadas subjetivas en decisiones fundamentadas.

\section{Conceptos Fundamentales}

\begin{definicion}{Experimento, Espacio Muestral y Eventos}
\begin{itemize}
    \item \textbf{Experimento Aleatorio:} Proceso mediante el cual se genera una observación o medición, cuyo resultado exacto no se puede predecir con certeza de antemano (ej. auditar una factura para ver si tiene o no error tributario).
    \item \textbf{Espacio Muestral ($S$):} Conjunto de todos los resultados elementales posibles de un experimento aleatorio.
    \item \textbf{Punto Muestral ($e_i$):} Cada uno de los resultados elementales e indivisibles que conforman el espacio muestral.
    \item \textbf{Evento o Suceso ($A$):} Cualquier subconjunto del espacio muestral ($A \subseteq S$).
\end{itemize}
\end{definicion}

\section{Enfoques Conceptuales de la Probabilidad}

\subsection{1. Enfoque Clásico o \textit{A Priori} (Laplace)}
Se aplica cuando todos los resultados elementales del espacio muestral son \textbf{equiprobables} (tienen la misma posibilidad de ocurrir).
\begin{equation}
    P(A) = \frac{N(A)}{N(S)} = \frac{\text{Número de resultados favorables al evento } A}{\text{Número total de resultados posibles en } S}
\end{equation}

\subsection{2. Enfoque de Frecuencia Relativa o Empírico}
Se basa en la observación histórica de la proporción de veces que ocurre un evento al repetir el experimento un número grande $n$ de veces bajo condiciones similares:
\begin{equation}
    P(A) \approx \frac{n(A)}{n} = \frac{\text{Número de veces que ocurrió } A}{\text{Número total de ensayos realizados}}
\end{equation}
Por la \textbf{Ley de los Grandes Números}, a medida que el número de ensayos $n$ tiende a infinito, la frecuencia relativa converge a la verdadera probabilidad teórica.

\subsection{3. Enfoque Subjetivo o Personalista}
Se utiliza cuando no se dispone de datos históricos repetitivos ni de resultados equiprobables. Representa el \textbf{grado de creencia o confianza} que un decisor u experto asigna a la ocurrencia de un evento, con base en su experiencia, información cualitativa y juicio profesional.

\section{Axiomas de Kolmogorov y Propiedades Fundamentales}
Para cualquier evento $A$ definido en un espacio muestral $S$:
\begin{enumerate}
    \item \textbf{Axioma de No Negatividad:} $0 \le P(A) \le 1$ para todo evento $A$.
    \item \textbf{Axioma del Suceso Seguro:} $P(S) = 1$.
    \item \textbf{Suceso Imposible:} $P(\emptyset) = 0$.
    \item \textbf{Regla del Complemento:} La suma de la probabilidad de que ocurra $A$ más la de que no ocurra ($A'$) es igual a 1:
    \begin{equation}
        P(A) + P(A') = 1 \implies P(A') = 1 - P(A)
    \end{equation}
\end{enumerate}

\section{Reglas de la Adición de Probabilidades}

\subsection{Regla Especial de la Adición (Eventos Mutuamente Excluyentes)}
Dos eventos $A$ y $B$ son \textbf{mutuamente excluyentes} (o disjuntos) si no pueden ocurrir simultáneamente ($A \cap B = \emptyset$).
\begin{equation}
    P(A \cup B) = P(A) + P(B)
\end{equation}

\subsection{Regla General de la Adición (Eventos No Excluyentes)}
Si los eventos $A$ y $B$ pueden ocurrir juntos ($A \cap B \neq \emptyset$):
\begin{equation}
    P(A \cup B) = P(A) + P(B) - P(A \cap B)
\end{equation}
donde $P(A \cap B)$ representa la probabilidad conjunta de que ocurran ambos eventos.

\begin{ejemplo}{Auditoría de Cumplimiento Tributario (Basado en Anderson)}
Una firma auditora en Santa Cruz revisa los registros fiscales de 200 medianas empresas. Se encontró que:
\begin{itemize}
    \item El $40\%$ de las empresas presenta inconsistencias en el Libro de Compras IVA (Evento $A$).
    \item El $25\%$ presenta inconsistencias en el Impuesto a las Utilidades IUE (Evento $B$).
    \item El $15\%$ presenta inconsistencias en \textbf{ambos} tributos ($A \cap B$).
\end{itemize}
Si se selecciona una empresa al azar, ¿cuál es la probabilidad de que tenga problemas tributarios en el IVA o en el IUE (o en ambos)?

\tcbline
\textbf{Solución:}
Identificamos las probabilidades dadas:
\[
P(A) = 0.40, \quad P(B) = 0.25, \quad P(A \cap B) = 0.15
\]
Aplicando la Regla General de la Adición:
\[
P(A \cup B) = P(A) + P(B) - P(A \cap B) = 0.40 + 0.25 - 0.15 = 0.50 \implies 50\%
\]
\textbf{Interpretación:} La probabilidad de que una empresa presente al menos una observación fiscal en IVA o IUE es del 50\%.
\end{ejemplo}

\section{Probabilidad Condicional e Independencia}

\subsection{Definición de Probabilidad Condicional}
\begin{definicion}{Probabilidad Condicional}
La probabilidad condicional del evento $A$, dado que se sabe con certeza que ya ocurrió el evento $B$ (con $P(B) > 0$), se define formalmente como:
\begin{equation}
    P(A|B) = \frac{P(A \cap B)}{P(B)}
\end{equation}
De manera análoga, la probabilidad de $B$ dado $A$ es:
\begin{equation}
    P(B|A) = \frac{P(A \cap B)}{P(A)}, \quad \text{con } P(A) > 0
\end{equation}
\end{definicion}

\subsection{Independencia Estadística de Eventos}
\begin{definicion}{Eventos Independientes}
Dos eventos $A$ y $B$ son estadísticamente independientes si la ocurrencia de uno no afecta en absoluto la probabilidad de ocurrencia del otro:
\begin{equation}
    P(A|B) = P(A) \quad \text{o equivalentemente} \quad P(B|A) = P(B)
\end{equation}
Si esta igualdad no se cumple, se dice que los eventos son \textbf{dependientes}.
\end{definicion}

\section{Reglas de la Multiplicación}

\subsection{Regla Especial de la Multiplicación (Eventos Independientes)}
Si $A$ y $B$ son independientes:
\begin{equation}
    P(A \cap B) = P(A) \cdot P(B)
\end{equation}

\subsection{Regla General de la Multiplicación (Eventos Dependientes)}
Para cualquier par de eventos:
\begin{equation}
    P(A \cap B) = P(A) \cdot P(B|A) = P(B) \cdot P(A|B)
\end{equation}

\section{Tablas de Contingencia y Probabilidad Conjunta}
Una tabla de contingencia clasifica una muestra o población respecto a dos variables categóricas simultáneas.

\begin{ejemplo}{Riesgo Crediticio y Tamaño de Empresa (Basado en Lind)}
El departamento de riesgo crediticio de un banco comercial analizó el historial de 500 créditos corporativos otorgados durante el último año, arrojando los siguientes datos:

\begin{center}
\begin{tabular}{lccc}
\toprule
\textbf{Categoría de Crédito} & \textbf{Microempresa ($M$)} & \textbf{Gran Empresa ($G$)} & \textbf{Total} \\
\midrule
Al Día / Cumplido ($C$) & 260 & 190 & \textbf{450} \\
En Mora / Incumplido ($D$) & 40 & 10 & \textbf{50} \\
\midrule
\textbf{Total} & \textbf{300} & \textbf{200} & \textbf{500} \\
\bottomrule
\end{tabular}
\end{center}

Si se elige un crédito al azar, determine:
\begin{enumerate}[label=\alph*)]
    \item La probabilidad de que el crédito corresponda a una Microempresa: $P(M)$.
    \item La probabilidad conjunta de que sea de una Microempresa y esté en Mora: $P(M \cap D)$.
    \item La probabilidad de que esté en Mora, dado que es de una Microempresa: $P(D|M)$.
    \item La probabilidad de que esté en Mora, dado que es de una Gran Empresa: $P(D|G)$.
    \item ¿Es el estado del crédito independiente del tamaño de la empresa?
\end{enumerate}

\tcbline
\textbf{Solución:}
\begin{enumerate}[label=\alph*)]
    \item $P(M) = \frac{300}{500} = 0.60 \implies 60\%$.
    \item $P(M \cap D) = \frac{40}{500} = 0.08 \implies 8\%$.
    \item $P(D|M) = \frac{P(M \cap D)}{P(M)} = \frac{40/500}{300/500} = \frac{40}{300} = 0.1333 \implies 13.33\%$.
    \item $P(D|G) = \frac{P(G \cap D)}{P(G)} = \frac{10/500}{200/500} = \frac{10}{200} = 0.0500 \implies 5.00\%$.
    \item \textbf{Evaluación de Independencia:}
    La probabilidad global de mora es $P(D) = \frac{50}{500} = 0.10$ ($10\%$).\\
    Dado que $P(D|M) = 13.33\% \neq P(D) = 10\%$, el tamaño de la empresa y la condición de mora \textbf{son eventos dependientes}. Las microempresas presentan una tasa de mora significativamente superior a la de las grandes corporaciones.
\end{enumerate}
\end{ejemplo}

\section{Teorema de la Probabilidad Total y Teorema de Bayes}

\subsection{Teorema de la Probabilidad Total}
Si los eventos $A_1, A_2, \dots, A_k$ forman una \textbf{partición} del espacio muestral $S$ (son mutuamente excluyentes dos a dos y su unión conforma todo $S$, con $P(A_i) > 0$), entonces para cualquier evento arbitrario $B \subseteq S$:
\begin{equation}
    P(B) = \sum_{i=1}^k P(A_i) \cdot P(B|A_i) = P(A_1)P(B|A_1) + P(A_2)P(B|A_2) + \dots + P(A_k)P(B|A_k)
\end{equation}

\subsection{Teorema de Bayes}
\begin{definicion}{Teorema de Bayes}
Bajo las mismas condiciones de partición, la probabilidad condicional \textit{a posteriori} de un evento causante $A_j$, dado que ya se ha observado la evidencia $B$, está dada por:
\begin{equation}
    P(A_j|B) = \frac{P(A_j) \cdot P(B|A_j)}{\sum_{i=1}^k P(A_i) \cdot P(B|A_i)}
\end{equation}
\end{definicion}

\begin{ejemplo}{Detección de Fraude en Transacciones Electrónicas (Basado en Anderson)}
Un sistema contable centralizado procesa órdenes de pago de tres departamentos: Ventas ($A_1$), Adquisiciones ($A_2$) y Logística ($A_3$).
\begin{itemize}
    \item Ventas genera el $50\%$ de las órdenes ($P(A_1) = 0.50$).
    \item Adquisiciones genera el $30\%$ ($P(A_2) = 0.30$).
    \item Logística genera el $20\%$ ($P(A_3) = 0.20$).
\end{itemize}
La auditoría interna ha detectado que la probabilidad de que una transacción presente irregularidades contables o sospecha de fraude (Evento $F$) según su departamento de origen es:
\[
P(F|A_1) = 0.01 \quad (1\%), \qquad P(F|A_2) = 0.04 \quad (4\%), \qquad P(F|A_3) = 0.02 \quad (2\%)
\]
Durante un control de rutina, el sistema de alertas tempranas detecta una transacción con irregularidades graves. ¿Cuál es la probabilidad de que dicha transacción fraudulenta provenga del departamento de Adquisiciones ($A_2$)?

\tcbline
\textbf{Solución Paso a Paso:}
\begin{enumerate}
    \item \textbf{Cálculo de la probabilidad global de irregularidad $P(F)$ (Probabilidad Total):}
    \begin{align*}
        P(F) &= P(A_1)P(F|A_1) + P(A_2)P(F|A_2) + P(A_3)P(F|A_3) \\
        &= (0.50)(0.01) + (0.30)(0.04) + (0.20)(0.02) \\
        &= 0.005 + 0.012 + 0.004 = 0.021 \quad (2.1\%)
    \end{align*}
    \item \textbf{Aplicación del Teorema de Bayes para calcular $P(A_2|F)$:}
    \[
    P(A_2|F) = \frac{P(A_2)P(F|A_2)}{P(F)} = \frac{(0.30)(0.04)}{0.021} = \frac{0.012}{0.021} \approx 0.5714 \implies 57.14\%
    \]
\end{enumerate}
\textbf{Interpretación Contable:} Aunque el departamento de Adquisiciones solo origina el 30\% del volumen total de transacciones (probabilidad a priori), si se detecta una irregularidad, la probabilidad de que provenga de Adquisiciones se eleva drásticamente al \textbf{57.14\%} (probabilidad a posteriori). La auditoría forense debe priorizar la revisión en dicha área.
\end{ejemplo}

\section{Formulario Resumen del Capítulo 2}
\begin{formulario}[Resumen de Fórmulas — Teoría de Probabilidades]
\begin{itemize}
    \item \textbf{Probabilidad Clásica:} $P(A) = \dfrac{N(A)}{N(S)}$
    \item \textbf{Regla del Complemento:} $P(A') = 1 - P(A)$
    \item \textbf{Regla General de la Adición:} $P(A \cup B) = P(A) + P(B) - P(A \cap B)$
    \item \textbf{Probabilidad Condicional:} $P(A|B) = \dfrac{P(A \cap B)}{P(B)}$
    \item \textbf{Regla de Multiplicación (General):} $P(A \cap B) = P(A) \cdot P(B|A)$
    \item \textbf{Regla de Multiplicación (Independientes):} $P(A \cap B) = P(A) \cdot P(B)$
    \item \textbf{Teorema de Bayes:} $P(A_j|B) = \dfrac{P(A_j) \cdot P(B|A_j)}{\sum_{i=1}^k P(A_i) \cdot P(B|A_i)}$
\end{itemize}
\end{formulario}

\section{Ejercicios Propuestos de la Unidad}

\subsection*{Nivel 1: Conceptos y Reglas de Probabilidad}
\begin{enumerate}
    \item Enuncie los tres axiomas fundamentales de probabilidad de Kolmogorov y explique su significado en términos de certeza e imposibilidad.
    \item Sean $A$ y $B$ dos eventos asociados a un experimento aleatorio tales que $P(A) = 0.50$, $P(B) = 0.30$ y $P(A \cap B) = 0.15$.
    \begin{enumerate}[label=\alph*)]
        \item Calcule la probabilidad de la unión $P(A \cup B)$.
        \item Calcule la probabilidad condicional $P(A|B)$.
        \item Determine si los eventos $A$ y $B$ son estadísticamente independientes. Justifique formalmente.
    \end{enumerate}
    \item Si $P(E_1) = 0.40$, $P(E_2) = 0.35$ y $E_1$ y $E_2$ son mutuamente excluyentes, calcule:
    \begin{enumerate}[label=\alph*)]
        \item $P(E_1 \cup E_2)$
        \item $P(E_1 \cap E_2)$
        \item $P(E_1|E_2)$
    \end{enumerate}
\end{enumerate}

\subsection*{Nivel 2: Tablas de Contingencia y Probabilidad Condicional}
\begin{enumerate}[resume]
    \item Una firma de auditoría clasificó una muestra de 200 empresas según su tamaño (Pequeña, Mediana, Grande) y su cumplimiento en la presentación de estados financieros a tiempo (A tiempo, Con retraso):
    \begin{center}
    \small
    \begin{tabular}{l|cc|c}
    \toprule
    \textbf{Tamaño de Empresa} & \textbf{A Tiempo ($T$)} & \textbf{Con Retraso ($R$)} & \textbf{Total} \\
    \midrule
    Pequeña ($P$) & 60 & 40 & 100 \\
    Mediana ($M$) & 45 & 15 & 60 \\
    Grande ($G$) & 35 & 5 & 40 \\
    \midrule
    \textbf{Total} & \textbf{140} & \textbf{60} & \textbf{200} \\
    \bottomrule
    \end{tabular}
    \end{center}
    Si se selecciona una empresa al azar:
    \begin{enumerate}[label=\alph*)]
        \item Calcule la probabilidad de que haya presentado sus estados financieros con retraso.
        \item Calcule la probabilidad condicional de que presente con retraso dado que es una empresa pequeña: $P(R|P)$.
        \item Calcule la probabilidad condicional de que sea una empresa grande dado que presentó a tiempo: $P(G|T)$.
        \item ¿Es el cumplimiento tributario independiente del tamaño de la empresa? Justifique.
    \end{enumerate}
\end{enumerate}

\subsection*{Nivel 3: Teorema de Bayes y Toma de Decisiones}
\begin{enumerate}[resume]
    \item En una compañía comercializadora de granos en Montero, el 70\% de los clientes paga con transferencia bancaria electrónica y el 30\% con cheque. De los pagos efectuados con cheque, el 10\% son rechazados por falta de fondos o firmas disconformes. De los pagos por transferencia, solo el 1\% presenta inconsistencias operativas. Si el departamento de cobranzas detecta un pago con problemas:
    \begin{enumerate}[label=\alph*)]
        \item ¿Cuál es la probabilidad global de que un pago elegido al azar tenga problemas?
        \item ¿Cuál es la probabilidad de que el pago observado haya sido efectuado mediante cheque?
    \end{enumerate}
    \item Una entidad bancaria clasifica a sus solicitantes de crédito hipotecario en tres categorías de riesgo: Bajo ($60\%$), Medio ($30\%$) y Alto ($10\%$). La probabilidad histórica de entrar en mora prolongada es del 1\% para la categoría de riesgo bajo, 5\% para riesgo medio y 20\% para riesgo alto. Si un cliente incurre en mora prolongada, determine la probabilidad a posteriori de que pertenezca a la categoría de riesgo alto.
\end{enumerate}


% Capítulo 3: Distribuciones Discretas
\chapter{Distribuciones de Probabilidad para Variables Aleatorias Discretas}
\label{cap:distribuciones_discretas}

\section*{Objetivos de Aprendizaje de la Unidad}
Al finalizar el estudio de esta unidad, el estudiante será capaz de:
\begin{enumerate}
    \item Definir y calcular la \textbf{función de probabilidad}, el \textbf{valor esperado} (media $\mu$) y la \textbf{varianza} ($\sigma^2$) de una variable aleatoria discreta.
    \item Identificar las condiciones de aplicación y calcular probabilidades con la \textbf{Distribución Binomial} en auditoría de cumplimiento y control de calidad.
    \item Modelar procesos de muestreo sin reemplazo en poblaciones finitas utilizando la \textbf{Distribución Hipergeométrica}.
    \item Aplicar la \textbf{Distribución de Poisson} para modelar la tasa de llegadas de clientes a bancos y eventos raros en sistemas de información contable.
    \item Utilizar la \textbf{aproximación de Poisson a probabilidades binomiales} cuando el tamaño de muestra es grande y la probabilidad de éxito es pequeña.
\end{enumerate}

\section{Variables Aleatorias Discretas y sus Distribuciones}

\begin{definicion}{Variable Aleatoria Discreta}
Una \textbf{variable aleatoria} $X$ es una función que asigna un valor numérico real a cada resultado del espacio muestral de un experimento aleatorio. Se denomina \textbf{discreta} si el conjunto de sus valores posibles es finito o infinito numerable (números enteros $0, 1, 2, 3, \dots$).
\end{definicion}

\subsection{Propiedades de la Función de Probabilidad $P(X=x) = f(x)$}
\begin{enumerate}
    \item $f(x) \ge 0$ para todo valor $x$.
    \item $\sum_{\text{todos los } x} f(x) = 1$.
\end{enumerate}

\subsection{Valor Esperado y Varianza de una Variable Discreta}
\begin{definicion}{Valor Esperado, Varianza y Desviación Estándar}
\begin{itemize}
    \item \textbf{Valor Esperado (Media Poblacional $\mu$):}
    \begin{equation}
        E(X) = \mu = \sum x \cdot f(x)
    \end{equation}
    \item \textbf{Varianza ($\sigma^2$):}
    \begin{equation}
        Var(X) = \sigma^2 = E[(X - \mu)^2] = \sum (x - \mu)^2 f(x) = \left[ \sum x^2 f(x) \right] - \mu^2
    \end{equation}
    \item \textbf{Desviación Estándar ($\sigma$):}
    \begin{equation}
        \sigma = \sqrt{\sigma^2}
    \end{equation}
\end{itemize}
\end{definicion}

\section{Distribución de Bernoulli}
Un ensayo de Bernoulli es un experimento que solo admite dos resultados posibles y mutuamente excluyentes: \textbf{Éxito} ($X=1$) con probabilidad $p$, o \textbf{Fracaso} ($X=0$) con probabilidad $q = 1-p$.
\begin{equation}
    E(X) = p, \qquad Var(X) = p(1-p) = pq
\end{equation}

\section{Distribución Binomial}
\begin{definicion}{Proceso y Distribución Binomial}
Un experimento binomial consiste en $n$ repeticiones independientes de ensayos de Bernoulli, donde:
\begin{enumerate}
    \item El número de ensayos $n$ es fijo.
    \item Cada ensayo solo tiene dos resultados posibles: Éxito ($p$) y Fracaso ($q = 1-p$).
    \item La probabilidad de éxito $p$ permanece constante de un ensayo a otro.
    \item Los ensayos son estadísticamente independientes entre sí.
\end{enumerate}
La función de probabilidad del número de éxitos $X$ en $n$ ensayos es:
\begin{equation}
    P(X = x) = \binom{n}{x} p^x (1-p)^{n-x} = \frac{n!}{x!(n-x)!} p^x q^{n-x}, \quad \text{para } x = 0, 1, 2, \dots, n
\end{equation}
\end{definicion}

\subsection{Parámetros de la Distribución Binomial}
\begin{itemize}
    \item \textbf{Media o Valor Esperado:} $\mu = E(X) = n \cdot p$
    \item \textbf{Varianza:} $\sigma^2 = n \cdot p \cdot q = n \cdot p \cdot (1-p)$
    \item \textbf{Desviación Estándar:} $\sigma = \sqrt{n \cdot p \cdot q}$
\end{itemize}

\begin{ejemplo}{Auditoría de Facturación Electrónica (Basado en Anderson)}
Una auditoría preliminar en una cadena de supermercados determinó que el $10\%$ ($p = 0.10$) de las facturas electrónicas emitidas contienen algún error de registro contable o glosa incompleta. El auditor tributario selecciona al azar una muestra de $n = 6$ facturas para su revisión detallada. Calcule:
\begin{enumerate}[label=\alph*)]
    \item La probabilidad de que exactamente 2 facturas contengan errores.
    \item La probabilidad de que ninguna factura contenga errores.
    \item La probabilidad de que al menos 2 facturas contengan errores.
    \item El número esperado y la desviación estándar de facturas erróneas.
\end{enumerate}

\tcbline
\textbf{Solución Paso a Paso:}
Identificamos los parámetros: $n = 6$, $p = 0.10$, $q = 1 - 0.10 = 0.90$.
\begin{enumerate}[label=\alph*)]
    \item \textbf{Exactamente 2 facturas erróneas ($x = 2$):}
    \[
    P(X = 2) = \binom{6}{2} (0.10)^2 (0.90)^{6-2} = 15 \cdot (0.01) \cdot (0.90)^4 = 15 \cdot 0.01 \cdot 0.6561 = 0.0984 \implies 9.84\%
    \]
    \item \textbf{Ninguna factura con errores ($x = 0$):}
    \[
    P(X = 0) = \binom{6}{0} (0.10)^0 (0.90)^6 = 1 \cdot 1 \cdot 0.531441 \approx 0.5314 \implies 53.14\%
    \]
    \item \textbf{Al menos 2 facturas con errores ($P(X \ge 2)$):}
    Usamos la regla del complemento para simplificar el cálculo:
    \[
    P(X \ge 2) = 1 - [P(X = 0) + P(X = 1)]
    \]
    Calculamos $P(X = 1)$:
    \[
    P(X = 1) = \binom{6}{1} (0.10)^1 (0.90)^5 = 6 \cdot 0.10 \cdot 0.59049 = 0.3543
    \]
    Sustituyendo:
    \[
    P(X \ge 2) = 1 - (0.5314 + 0.3543) = 1 - 0.8857 = 0.1143 \implies 11.43\%
    \]
    \item \textbf{Media y Desviación Estándar:}
    \begin{align*}
        \mu &= n \cdot p = 6 \cdot 0.10 = 0.60 \text{ facturas con error en promedio.} \\
        \sigma &= \sqrt{n \cdot p \cdot q} = \sqrt{6 \cdot 0.10 \cdot 0.90} = \sqrt{0.54} \approx 0.7348 \text{ facturas.}
    \end{align*}
\end{enumerate}
\end{ejemplo}

\section{Distribución Hipergeométrica}
\begin{definicion}{Distribución Hipergeométrica (Muestreo Sin Reemplazo)}
Cuando se extrae una muestra de tamaño $n$ \textbf{sin reemplazo} a partir de una población finita de tamaño $N$, donde existen $T$ elementos clasificados como éxitos y $N-T$ clasificados como fracasos, la probabilidad de obtener exactamente $x$ éxitos en la muestra es:
\begin{equation}
    P(X = x) = \frac{\binom{T}{x} \binom{N-T}{n-x}}{\binom{N}{n}}
\end{equation}
para $\max(0, n - (N-T)) \le x \le \min(n, T)$.
\end{definicion}

\subsection{Parámetros de la Hipergeométrica}
\begin{itemize}
    \item \textbf{Media:} $\mu = n \left( \dfrac{T}{N} \right)$
    \item \textbf{Varianza:} $\sigma^2 = n \left( \dfrac{T}{N} \right) \left( 1 - \dfrac{T}{N} \right) \left( \dfrac{N - n}{N - 1} \right)$
\end{itemize}
donde el término $\frac{N-n}{N-1}$ es el \textbf{Factor de Corrección por Población Finita (FCPF)}.

\begin{ejemplo}{Auditoría de Contratos Públicos (Basado en Lind)}
El departamento de auditoría interna de una entidad pública revisa un expediente de $N = 20$ contratos de compras adjudicados el mes anterior. Se sabe que $T = 5$ de estos contratos presentan irregularidades en la documentación de respaldo. El auditor jefe selecciona al azar una muestra de $n = 4$ contratos para fiscalización presencial sin reemplazo.
\begin{enumerate}[label=\alph*)]
    \item ¿Cuál es la probabilidad de que la muestra contenga exactamente 2 contratos con irregularidades?
    \item ¿Cuál es la probabilidad de no encontrar ningún contrato irregular en la muestra?
\end{enumerate}

\tcbline
\textbf{Solución:}
Datos: $N = 20$, $T = 5$ irregulares, $N - T = 15$ conformes, $n = 4$.
\begin{enumerate}[label=\alph*)]
    \item \textbf{Exactamente 2 irregulares ($x = 2$):}
    \[
    P(X = 2) = \frac{\binom{5}{2} \binom{15}{4-2}}{\binom{20}{4}} = \frac{\binom{5}{2} \binom{15}{2}}{\binom{20}{4}}
    \]
    Calculamos cada combinatorio:
    \begin{align*}
        \binom{5}{2} &= 10, \qquad \binom{15}{2} = \frac{15 \cdot 14}{2} = 105, \qquad \binom{20}{4} = \frac{20 \cdot 19 \cdot 18 \cdot 17}{4 \cdot 3 \cdot 2 \cdot 1} = 4{,}845 \\
        P(X = 2) &= \frac{10 \cdot 105}{4{,}845} = \frac{1{,}050}{4{,}845} \approx 0.2167 \implies 21.67\%
    \end{align*}
    \item \textbf{Ningún contrato irregular ($x = 0$):}
    \[
    P(X = 0) = \frac{\binom{5}{0} \binom{15}{4}}{\binom{20}{4}} = \frac{1 \cdot 1{,}365}{4{,}845} \approx 0.2817 \implies 28.17\%
    \]
\end{enumerate}
\end{ejemplo}

\section{Distribución de Poisson}
\begin{definicion}{Distribución de Poisson}
Se aplica a experimentos donde se cuenta el número de veces $X$ que ocurre un evento en un \textbf{intervalo continuo} determinado de tiempo, área, volumen o distancia.
\begin{equation}
    P(X = x) = \frac{\lambda^x e^{-\lambda}}{x!}, \quad \text{para } x = 0, 1, 2, 3, \dots
\end{equation}
donde $\lambda$ ($\lambda > 0$) es la tasa promedio de ocurrencias por intervalo y $e \approx 2.71828$.
\end{definicion}

\subsection{Parámetros de la Distribución de Poisson}
\begin{itemize}
    \item \textbf{Media:} $\mu = E(X) = \lambda$
    \item \textbf{Varianza:} $\sigma^2 = Var(X) = \lambda$
    \item \textbf{Desviación Estándar:} $\sigma = \sqrt{\lambda}$
\end{itemize}

\begin{ejemplo}{Atención a Clientes en Cajas Bancarias (Basado en Anderson)}
En una agencia bancaria del centro de Santa Cruz, los clientes llegan a la ventanilla de depósitos a razón promedio de $\lambda = 3$ clientes cada 10 minutos durante las horas pico. Calcule:
\begin{enumerate}[label=\alph*)]
    \item La probabilidad de que lleguen exactamente 4 clientes en un lapso de 10 minutos.
    \item La probabilidad de que no llegue ningún cliente en 10 minutos.
    \item La probabilidad de que lleguen más de 2 clientes en 10 minutos.
    \item La probabilidad de que lleguen exactamente 2 clientes en un periodo de 20 minutos.
\end{enumerate}

\tcbline
\textbf{Solución Paso a Paso:}
\begin{enumerate}[label=\alph*)]
    \item \textbf{Exactamente 4 clientes en 10 min ($\lambda = 3, x = 4$):}
    \[
    P(X = 4) = \frac{3^4 e^{-3}}{4!} = \frac{81 \cdot 0.049787}{24} = \frac{4.03275}{24} = 0.1680 \implies 16.80\%
    \]
    \item \textbf{Ningún cliente en 10 min ($x = 0$):}
    \[
    P(X = 0) = \frac{3^0 e^{-3}}{0!} = \frac{1 \cdot 0.049787}{1} \approx 0.0498 \implies 4.98\%
    \]
    \item \textbf{Más de 2 clientes en 10 min ($P(X > 2)$):}
    \[
    P(X > 2) = 1 - [P(X = 0) + P(X = 1) + P(X = 2)]
    \]
    Calculamos $P(X = 1)$ y $P(X = 2)$:
    \begin{align*}
        P(X = 1) &= \frac{3^1 e^{-3}}{1!} = 3(0.049787) = 0.1494 \\
        P(X = 2) &= \frac{3^2 e^{-3}}{2!} = \frac{9(0.049787)}{2} = 0.2240 \\
        P(X > 2) &= 1 - (0.0498 + 0.1494 + 0.2240) = 1 - 0.4232 = 0.5768 \implies 57.68\%
    \end{align*}
    \item \textbf{Llegadas en un periodo de 20 minutos:}
    Al duplicar el intervalo de tiempo, la tasa esperada se ajusta proporcionalmente: $\lambda_{20\text{min}} = 3 \times 2 = 6$.
    \[
    P(X = 2) = \frac{6^2 e^{-6}}{2!} = \frac{36 \cdot 0.002479}{2} = 18 \cdot 0.002479 = 0.0446 \implies 4.46\%
    \]
\end{enumerate}
\end{ejemplo}

\section{Aproximación de Poisson a la Distribución Binomial}
Cuando en un experimento binomial el número de ensayos $n$ es muy grande ($n \ge 20$ o preferentemente $n \ge 100$) y la probabilidad de éxito $p$ es muy pequeña ($p \le 0.05$ o $np \le 10$), la distribución binomial converge a una distribución de Poisson con parámetro:
\begin{equation}
    \lambda = n \cdot p
\end{equation}

\begin{ejemplo}{Detección de Errores Raros en Transferencias Interbancarias}
Un sistema interbancario procesa $n = 2{,}000$ transferencias diarias. La probabilidad de que una transferencia individual sufra una falla técnica de acreditación es de $p = 0.0015$. Utilice la aproximación de Poisson para calcular la probabilidad de que se presenten exactamente 3 fallas en un día.

\tcbline
\textbf{Solución:}
Comprobamos condiciones: $n = 2{,}000 \ge 100$ y $p = 0.0015 \le 0.05$.
Calculamos la tasa media: $\lambda = n \cdot p = 2{,}000 \times 0.0015 = 3.0$.
\[
P(X = 3) = \frac{3^3 e^{-3}}{3!} = \frac{27 \cdot 0.049787}{6} = 0.2240 \implies 22.40\%
\]
\end{ejemplo}

\section{Formulario Resumen del Capítulo 3}
\begin{formulario}[Resumen de Fórmulas — Distribuciones Discretas]
\begin{itemize}
    \item \textbf{Valor Esperado Discreto:} $\mu = \sum x f(x), \quad \sigma^2 = \sum x^2 f(x) - \mu^2$
    \item \textbf{Binomial:} $P(X=x) = \dbinom{n}{x} p^x q^{n-x}, \quad \mu = np, \quad \sigma = \sqrt{npq}$
    \item \textbf{Hipergeométrica:} $P(X=x) = \dfrac{\dbinom{T}{x} \dbinom{N-T}{n-x}}{\dbinom{N}{n}}, \quad \mu = n\left(\dfrac{T}{N}\right)$
    \item \textbf{Poisson:} $P(X=x) = \dfrac{\lambda^x e^{-\lambda}}{x!}, \quad \mu = \lambda, \quad \sigma = \sqrt{\lambda}$
    \item \textbf{Aprox. Poisson a Binomial:} $\lambda = n \cdot p$ (si $n \ge 100$ y $p \le 0.05$)
\end{itemize}
\end{formulario}

\section{Ejercicios Propuestos de la Unidad}

\subsection*{Nivel 1: Identificación de Distribuciones y Cálculo Básico}
\begin{enumerate}
    \item Identifique y justifique qué modelo de distribución de probabilidad discreta (Binomial, Hipergeométrica o Poisson) es el más adecuado para modelar cada una de las siguientes situaciones:
    \begin{enumerate}[label=\alph*)]
        \item El número de clientes que ingresan a una agencia bancaria entre las 10:00 y las 11:00 de la mañana.
        \item El número de facturas con errores de retención en una muestra de 10 facturas extraídas con reemplazo de un talonario con 15\% de errores.
        \item El número de cheques sin fondos en una muestra de 5 cheques extraídos sin reemplazo de un lote de 20 cheques que contiene 3 sin fondos.
    \end{enumerate}
    \item Una variable aleatoria discreta $X$ sigue una distribución Binomial con $n = 10$ ensayos y probabilidad de éxito $p = 0.20$. Calcule:
    \begin{enumerate}[label=\alph*)]
        \item $P(X = 0)$
        \item $P(X = 2)$
        \item $P(X \le 2)$
        \item El valor esperado $E(X)$ y la desviación estándar $\sigma$.
    \end{enumerate}
\end{enumerate}

\subsection*{Nivel 2: Aplicaciones en Auditoría y Control de Calidad}
\begin{enumerate}[resume]
    \item Un banco sabe que el 5\% de los solicitantes de crédito personal presentan información laboral no comprobable o falsificada. Si un oficial de crédito revisa 8 solicitudes seleccionadas al azar:
    \begin{enumerate}[label=\alph*)]
        \item ¿Cuál es la probabilidad de encontrar exactamente 1 solicitud irregular?
        \item ¿Cuál es la probabilidad de encontrar al menos 1 solicitud irregular?
    \end{enumerate}
    \item Un lote de 25 computadoras para el laboratorio contable contiene 4 equipos con fallas en la memoria RAM. Si se eligen 5 computadoras al azar sin reemplazo para instalarlas en el departamento de auditoría:
    \begin{enumerate}[label=\alph*)]
        \item Calcule la probabilidad de que ningún equipo sea defectuoso.
        \item Calcule la probabilidad de encontrar exactamente 2 defectuosos.
        \item Calcule el número esperado de computadoras defectuosas en la muestra.
    \end{enumerate}
    \item En una central de atención telefónica de cobranzas, las consultas por deudas vencidas ingresan a una tasa media de $\lambda = 4$ llamadas por hora.
    \begin{enumerate}[label=\alph*)]
        \item Calcule la probabilidad de recibir exactamente 3 llamadas en una hora determinada.
        \item Calcule la probabilidad de recibir al menos 2 llamadas en una hora.
        \item Calcule la probabilidad de recibir exactamente 5 llamadas en un turno de 2 horas continuas ($\lambda = 8$).
    \end{enumerate}
\end{enumerate}

\subsection*{Nivel 3: Casos Integrales y Aproximaciones}
\begin{enumerate}[resume]
    \item Una compañía aseguradora en Santa Cruz tiene una cartera de $n = 5{,}000$ pólizas de seguro contra accidentes laborales para empresas constructoras. La probabilidad histórica de que un trabajador asegurado sufra un accidente grave que requiera indemnización máxima durante el año es de $p = 0.0006$.
    \begin{enumerate}[label=\alph*)]
        \item Justifique por qué es válido emplear la aproximación de Poisson a la distribución Binomial.
        \item Determine el parámetro $\lambda$ de la distribución de Poisson equivalente.
        \item Calcule la probabilidad de que la aseguradora deba pagar exactamente 2 indemnizaciones máximas en el año.
        \item Calcule la probabilidad de que deba pagar más de 3 indemnizaciones en el año.
    \end{enumerate}
\end{enumerate}


% Capítulo 4: Distribuciones Continuas
\chapter{Distribuciones de Probabilidad para Variables Aleatorias Continuas}
\label{cap:distribuciones_continuas}

\section*{Objetivos de Aprendizaje de la Unidad}
Al finalizar el estudio de esta unidad, el estudiante será capaz de:
\begin{enumerate}
    \item Comprender la naturaleza de las variables aleatorias continuas y la interpretación del área bajo la \textbf{función de densidad de probabilidad} $f(x)$.
    \item Dominar la \textbf{Distribución Normal} y el proceso de estandarización con la variable $Z$.
    \item Calcular probabilidades directas e inversas (percentiles y valores críticos) en problemas de costos, salarios y finanzas corporativas.
    \item Aplicar la \textbf{Distribución Exponencial} para modelar tiempos de espera y liquidación de trámites.
    \item Realizar la \textbf{aproximación Normal a la distribución Binomial} utilizando adecuadamente el factor de corrección por continuidad.
\end{enumerate}

\section{Variables Aleatorias Continuas y Funciones de Densidad}

\begin{definicion}{Variable Aleatoria Continua y Función de Densidad}
Una variable aleatoria $X$ es \textbf{continua} si puede tomar cualquier valor numérico dentro de un intervalo continuo de la recta real (ej. ingresos mensuales, costos de producción, tiempo de auditoría, volumen de ventas). Su comportamiento probabilístico se describe mediante una \textbf{Función de Densidad de Probabilidad (FDP)} $f(x)$ que cumple:
\begin{enumerate}
    \item $f(x) \ge 0$ para todo $x \in \mathbb{R}$.
    \item $\int_{-\infty}^{\infty} f(x) \, dx = 1$ (el área total bajo la curva es igual a 1).
    \item La probabilidad de que $X$ se encuentre en el intervalo $[a, b]$ es el área bajo la curva entre $a$ y $b$:
    \begin{equation}
        P(a \le X \le b) = \int_a^b f(x) \, dx
    \end{equation}
\end{enumerate}
\end{definicion}

\begin{alerta}
Para cualquier variable aleatoria continua, la probabilidad de que tome un valor puntual exacto es siempre \textbf{cero}:
\[
P(X = c) = 0 \quad \text{para cualquier constante } c.
\]
En consecuencia, las desigualdades estrictas y no estrictas son equivalentes:
\[
P(a \le X \le b) = P(a < X \le b) = P(a \le X < b) = P(a < X < b)
\]
\end{alerta}

\section{La Distribución Normal o Gaussiana}
La distribución normal es la más importante y utilizada en estadística teórica y aplicada. Fue descrita por Carl Friedrich Gauss y modela una inmensa cantidad de fenómenos económicos, financieros y naturales.

\begin{definicion}{Función de Densidad Normal}
Una variable aleatoria continua $X$ con media $\mu$ ($-\infty < \mu < \infty$) y desviación estándar $\sigma$ ($\sigma > 0$) tiene una \textbf{Distribución Normal} ($X \sim N(\mu, \sigma^2)$) si su función de densidad es:
\begin{equation}
    f(x) = \frac{1}{\sigma \sqrt{2\pi}} e^{-\frac{1}{2}\left(\frac{x - \mu}{\sigma}\right)^2}, \quad -\infty < x < \infty
\end{equation}
donde $\pi \approx 3.14159$ y $e \approx 2.71828$.
\end{definicion}

\subsection{Propiedades Geométricas y Estadísticas de la Campana Normal}
\begin{enumerate}
    \item \textbf{Forma Acampanada y Simétrica:} La curva es perfectamente simétrica respecto a la media $\mu$. La media, la mediana y la moda coinciden exactamente en el centro ($\mu = \text{Mediana} = \text{Moda}$).
    \item \textbf{Asíntotas Horizontales:} Los extremos de la curva se extienden indefinidamente hacia $-\infty$ y $+\infty$ acercándose asintóticamente al eje horizontal sin llegar a tocarlo jamás.
    \item \textbf{Regla Empírica de Probabilidad:}
    \begin{itemize}
        \item $\mu \pm 1\sigma$ contiene aproximadamente el \textbf{68.26\%} del área total.
        \item $\mu \pm 2\sigma$ contiene aproximadamente el \textbf{95.44\%} del área total.
        \item $\mu \pm 3\sigma$ contiene aproximadamente el \textbf{99.74\%} del área total.
    \end{itemize}
\end{enumerate}

\section{La Distribución Normal Estándar ($Z$)}
Dado que existen infinitas distribuciones normales (según los valores de $\mu$ y $\sigma$), para calcular probabilidades se realiza una transformación algebraica lineal que estandariza cualquier variable normal $X$ a una variable \textbf{Normal Estándar $Z$} con media $\mu_Z = 0$ y desviación estándar $\sigma_Z = 1$ ($Z \sim N(0, 1)$).

\begin{definicion}{Estandarización $Z$}
El valor $Z$ (puntuación estándar) mide el número de desviaciones estándar que un valor particular $x$ se aleja de la media poblacional $\mu$:
\begin{equation}
    Z = \frac{X - \mu}{\sigma}
\end{equation}
\end{definicion}

\subsection{Cálculo de Probabilidades Directas}
\begin{ejemplo}{Análisis de Salarios en Empresas de Servicios (Basado en Anderson)}
En un estudio sobre el mercado laboral contable en Santa Cruz, se determinó que los salarios mensuales de los contadores junior tienen una distribución aproximadamente normal con una media de $\mu = 5{,}200$ Bs y una desviación estándar de $\sigma = 600$ Bs. Si se selecciona un contador junior al azar:
\begin{enumerate}[label=\alph*)]
    \item ¿Cuál es la probabilidad de que su salario mensual sea menor a $4{,}600$ Bs?
    \item ¿Cuál es la probabilidad de que su salario se encuentre entre $4{,}900$ Bs y $6{,}100$ Bs?
    \item ¿Cuál es la probabilidad de que su salario sea superior a $6{,}400$ Bs?
\end{enumerate}

\tcbline
\textbf{Solución Paso a Paso:}
Datos: $\mu = 5{,}200$, $\sigma = 600$.
\begin{enumerate}[label=\alph*)]
    \item \textbf{Salario menor a $4{,}600$ Bs ($P(X < 4{,}600)$):}
    Estandarizamos el valor $x = 4{,}600$:
    \[
    Z = \frac{4{,}600 - 5{,}200}{600} = \frac{-600}{600} = -1.00
    \]
    Buscando en la tabla de la normal estándar: $P(Z < -1.00) = 0.1587 \implies 15.87\%$.
    
    \item \textbf{Salario entre $4{,}900$ Bs y $6{,}100$ Bs ($P(4{,}900 < X < 6{,}100)$):}
    Estandarizamos ambos límites:
    \[
    Z_1 = \frac{4{,}900 - 5{,}200}{600} = \frac{-300}{600} = -0.50, \qquad Z_2 = \frac{6{,}100 - 5{,}200}{600} = \frac{900}{600} = 1.50
    \]
    Calculamos el área intermedia:
    \begin{align*}
        P(-0.50 < Z < 1.50) &= P(Z < 1.50) - P(Z < -0.50) \\
        &= 0.9332 - 0.3085 = 0.6247 \implies 62.47\%
    \end{align*}
    
    \item \textbf{Salario superior a $6{,}400$ Bs ($P(X > 6{,}400)$):}
    \[
    Z = \frac{6{,}400 - 5{,}200}{600} = \frac{1{,}200}{600} = 2.00
    \]
    \[
    P(Z > 2.00) = 1 - P(Z \le 2.00) = 1 - 0.9772 = 0.0228 \implies 2.28\%
    \]
\end{enumerate}
\end{ejemplo}

\subsection{Cálculo de Problemas Inversos (Percentiles y Puntos de Corte)}
En muchas decisiones financieras y de control de calidad se conoce la probabilidad acumulada requerida y se necesita encontrar el valor umbral $x$ correspondiente. Despejando $x$ de la fórmula de estandarización:
\begin{equation}
    x = \mu + Z \cdot \sigma
\end{equation}

\begin{ejemplo}{Fijación de Política de Bonos de Desempeño (Basado en Lind)}
La gerencia general de una corporación financiera decide otorgar un bono especial de productividad al $10\%$ superior de sus ejecutivos de cobranzas con mayor volumen recuperado al mes. Si el monto mensual recuperado sigue una distribución normal con $\mu = 85{,}000$ USD y $\sigma = 12{,}000$ USD, ¿cuál es el monto mínimo de recuperación que debe alcanzar un ejecutivo para ser acreedor al bono?

\tcbline
\textbf{Solución:}
El $10\%$ superior equivale a decir que el $90\%$ ($0.90$) de los ejecutivos recupera un monto inferior a dicho umbral ($P(X \le x) = 0.90$).
\begin{enumerate}
    \item Buscamos en la tabla $Z$ el valor que deja un área acumulada de $0.9000$:
    \[
    Z_{0.90} \approx 1.28
    \]
    \item Despejamos el valor $x$:
    \[
    x = \mu + Z \cdot \sigma = 85{,}000 + (1.28)(12{,}000) = 85{,}000 + 15{,}360 = 100{,}360 \text{ USD.}
    \]
\end{enumerate}
\textbf{Interpretación Financiera:} Para acceder al bono de productividad, el ejecutivo debe recuperar como mínimo \textbf{100,360 USD} en el mes.
\end{ejemplo}

\section{Aproximación Normal a la Distribución Binomial}
Cuando el número de ensayos $n$ en un experimento binomial es grande, el cálculo combinatorio se vuelve laborioso. El Teorema de De Moivre-Laplace establece que si se cumplen las condiciones:
\begin{equation}
    n \cdot p \ge 5 \quad \text{y} \quad n \cdot q \ge 5
\end{equation}
la variable binomial discreta $X$ puede aproximarse mediante una distribución normal continua con:
\begin{equation}
    \mu = n \cdot p, \qquad \sigma = \sqrt{n \cdot p \cdot q}
\end{equation}

\subsection{Factor de Corrección por Continuidad ($\pm 0.5$)}
Dado que se está aproximando una distribución discreta (escalonada) mediante una curva continua y suave, cada valor entero $x$ discreto se representa en el continuo como el intervalo $[x - 0.5, x + 0.5]$.
\begin{itemize}
    \item $P(X = k) \approx P(k - 0.5 \le X \le k + 0.5)$
    \item $P(X \le k) \approx P(X \le k + 0.5)$
    \item $P(X \ge k) \approx P(X \ge k - 0.5)$
    \item $P(a \le X \le b) \approx P(a - 0.5 \le X \le b + 0.5)$
\end{itemize}

\begin{ejemplo}{Auditoría Masiva de Cuentas por Cobrar (Basado en Anderson)}
Una tienda de departamentos tiene $n = 400$ cuentas de crédito activas. Históricamente se sabe que el $20\%$ ($p = 0.20$) de las cuentas caen en mora. Calcule la probabilidad de que entre 70 y 95 cuentas (ambas inclusive) se encuentren en mora.

\tcbline
\textbf{Solución:}
\begin{enumerate}
    \item \textbf{Verificación de supuestos:}
    \[
    np = 400(0.20) = 80 \ge 5, \qquad nq = 400(0.80) = 320 \ge 5
    \]
    \item \textbf{Parámetros de la Normal:}
    \[
    \mu = 80, \qquad \sigma = \sqrt{npq} = \sqrt{400(0.20)(0.80)} = \sqrt{64} = 8.0
    \]
    \item \textbf{Aplicación de la corrección por continuidad:}
    \[
    P(70 \le X \le 95) \approx P(69.5 \le X_{\text{normal}} \le 95.5)
    \]
    \item \textbf{Estandarización $Z$:}
    \[
    Z_1 = \frac{69.5 - 80}{8} = \frac{-10.5}{8} = -1.31, \qquad Z_2 = \frac{95.5 - 80}{8} = \frac{15.5}{8} = 1.94
    \]
    \[
    P(-1.31 \le Z \le 1.94) = P(Z \le 1.94) - P(Z \le -1.31) = 0.9738 - 0.0951 = 0.8787 \implies 87.87\%
    \]
\end{enumerate}
\end{ejemplo}

\section{La Distribución Exponencial}
\begin{definicion}{Distribución Exponencial}
Se utiliza frecuentemente para modelar el tiempo transcurrido hasta que ocurre un determinado suceso (ej. tiempo de atención a un contribuyente, vida útil de un equipo contable). Su función de densidad es:
\begin{equation}
    f(x) = \lambda e^{-\lambda x}, \quad \text{para } x \ge 0
\end{equation}
donde $\lambda > 0$ es la tasa media de ocurrencias por unidad de tiempo.
\end{definicion}

\subsection{Probabilidades Acumuladas y Parámetros}
\begin{itemize}
    \item \textbf{Probabilidad acumulada:} $P(X \le x) = 1 - e^{-\lambda x}$
    \item \textbf{Probabilidad de cola derecha:} $P(X > x) = e^{-\lambda x}$
    \item \textbf{Media y Desviación Estándar:} $\mu = \dfrac{1}{\lambda}, \quad \sigma = \dfrac{1}{\lambda}$
\end{itemize}

\begin{ejemplo}{Tiempo de Atención en Plataforma Tributaria}
El tiempo que tarda un funcionario del Servicio de Impuestos Nacionales en revisar un trámite de facturación sigue una distribución exponencial con un promedio de 15 minutos por trámite ($\mu = 15 \implies \lambda = 1/15 \approx 0.0667$ trámites por minuto).
\begin{enumerate}[label=\alph*)]
    \item ¿Cuál es la probabilidad de que un trámite tarde menos de 10 minutos?
    \item ¿Cuál es la probabilidad de que tarde más de 20 minutos?
\end{enumerate}

\tcbline
\textbf{Solución:}
\begin{enumerate}[label=\alph*)]
    \item $P(X \le 10) = 1 - e^{-(1/15)(10)} = 1 - e^{-0.6667} = 1 - 0.5134 = 0.4866 \implies 48.66\%$.
    \item $P(X > 20) = e^{-(1/15)(20)} = e^{-1.3333} \approx 0.2636 \implies 26.36\%$.
\end{enumerate}
\end{ejemplo}

\section{Formulario Resumen del Capítulo 4}
\begin{formulario}[Resumen de Fórmulas — Distribuciones Continuas]
\begin{itemize}
    \item \textbf{Estandarización Normal:} $Z = \dfrac{X - \mu}{\sigma}, \quad X = \mu + Z \cdot \sigma$
    \item \textbf{Aprox. Normal a Binomial:} $\mu = np, \quad \sigma = \sqrt{npq}$ (con corrección $\pm 0.5$)
    \item \textbf{Distribución Exponencial:} $P(X \le x) = 1 - e^{-\lambda x}, \quad P(X > x) = e^{-\lambda x}, \quad \mu = \dfrac{1}{\lambda}$
\end{itemize}
\end{formulario}

\section{Ejercicios Propuestos de la Unidad}

\subsection*{Nivel 1: Manejo de la Tabla Normal y Estandarización}
\begin{enumerate}
    \item Para una variable aleatoria normal estándar $Z \sim N(0, 1)$, utilice la tabla de distribución para calcular las siguientes probabilidades:
    \begin{enumerate}[label=\alph*)]
        \item $P(Z \le 1.65)$
        \item $P(Z > 2.10)$
        \item $P(-1.25 \le Z \le 1.85)$
        \item Encuentre el valor crítico $z_0$ tal que $P(Z \le z_0) = 0.95$.
        \item Encuentre el valor crítico $z_0$ tal que $P(-z_0 \le Z \le z_0) = 0.99$.
    \end{enumerate}
    \item Una variable aleatoria continua $X$ sigue una distribución Normal con media $\mu = 500$ y desviación estándar $\sigma = 50$. Calcule:
    \begin{enumerate}[label=\alph*)]
        \item $P(X \le 580)$
        \item $P(X > 430)$
        \item $P(460 \le X \le 540)$
        \item El percentil 90 ($P_{90}$).
    \end{enumerate}
\end{enumerate}

\subsection*{Nivel 2: Aplicaciones Económicas y Financieras}
\begin{enumerate}[resume]
    \item Los costos de honorarios por auditoría externa para pequeñas y medianas empresas en Santa Cruz tienen una distribución normal con media $\mu = 12{,}500$ Bs y desviación estándar $\sigma = 1{,}800$ Bs. Calcule la probabilidad de que una auditoría seleccionada al azar cueste:
    \begin{enumerate}[label=\alph*)]
        \item Menos de $10{,}000$ Bs.
        \item Más de $15{,}000$ Bs.
        \item Entre $11{,}000$ Bs y $14{,}000$ Bs.
        \item ¿Cuál es el costo máximo que pagará el 10\% de las empresas con auditorías más económicas (percentil 10)?
    \end{enumerate}
    \item El tiempo de atención por ventanilla a clientes corporativos en una entidad bancaria sigue una distribución exponencial con media de 12 minutos ($\lambda = 1/12 \approx 0.0833$ clientes por minuto).
    \begin{enumerate}[label=\alph*)]
        \item Calcule la probabilidad de que un cliente sea atendido en menos de 8 minutos.
        \item Calcule la probabilidad de que un cliente deba esperar más de 20 minutos.
        \item Calcule la probabilidad de que la atención dure entre 10 y 15 minutos.
    \end{enumerate}
\end{enumerate}

\subsection*{Nivel 3: Aproximación Normal a la Binomial y Casos de Auditoría}
\begin{enumerate}[resume]
    \item En una compañía de seguros, el 15\% de las pólizas de vehículos sufren siniestros en un año. Para una cartera de $n = 200$ asegurados independientes:
    \begin{enumerate}[label=\alph*)]
        \item Verifique si se cumplen los supuestos para aproximar la distribución Binomial mediante la Normal ($np \ge 5$ y $nq \ge 5$).
        \item Calcule la media $\mu$ y la desviación estándar $\sigma$ de la distribución normal aproximada.
        \item Utilizando el factor de corrección por continuidad, calcule la probabilidad de que ocurran exactamente 35 siniestros ($P(34.5 \le X \le 35.5)$).
        \item Calcule la probabilidad de que se presenten a lo sumo 25 siniestros.
    \end{enumerate}
    \item Un fondo mutuo de inversión en Bolivia reporta que los rendimientos anuales de sus instrumentos de renta fija se distribuyen normalmente con media $\mu = 7.5\%$ y desviación estándar $\sigma = 1.8\%$. Si un inversionista institucional coloca 1 millón de bolivianos en el fondo, ¿cuál es la probabilidad de que el rendimiento anual sea inferior al 4.0\% (pérdida en términos reales considerando la inflación)?
\end{enumerate}


% Capítulo 5: Muestreo e Intervalos de Confianza
\chapter{Distribuciones Muestrales e Intervalos de Confianza para la Media}
\label{cap:muestreo_intervalos}

\section*{Objetivos de Aprendizaje de la Unidad}
Al finalizar el estudio de esta unidad, el estudiante será capaz de:
\begin{enumerate}
    \item Comprender la fundamentación teórica del \textbf{muestreo probabilístico} en auditoría de estados financieros y control de gestión.
    \item Aplicar el \textbf{Teorema del Límite Central (TLC)} para determinar la distribución probabilística de la media muestral $\bar{X}$.
    \item Construir e interpretar \textbf{intervalos de confianza} para la media poblacional $\mu$ utilizando la distribución Normal ($Z$) y la distribución $t$ de Student.
    \item Calcular con rigor técnico el \textbf{tamaño óptimo de muestra ($n$)} requerido para auditorías y estudios económicos, controlando el error muestral y el nivel de confianza.
\end{enumerate}

\section{Fundamentos de la Inferencia Estadística en Auditoría}
En la práctica profesional de la contaduría pública y la auditoría, examinar la totalidad de los comprobantes contables de una entidad económica (censo) suele ser inviable debido a restricciones de tiempo, costo y operatividad. Las \textbf{Normas Internacionales de Auditoría (NIA 530 - Muestreo de Auditoría)} exigen que el auditor aplique procedimientos estadísticos rigurosos para extraer conclusiones válidas sobre la población total a partir de una muestra representativa.

\begin{definicion}{Parámetros vs. Estadísticos}
\begin{itemize}
    \item \textbf{Población ($N$):} Conjunto completo de todos los elementos, facturas, transacciones o cuentas que son objeto de estudio.
    \item \textbf{Muestra ($n$):} Subconjunto representativo seleccionado de la población ($n < N$).
    \item \textbf{Parámetro (Letras Griegas):} Medida cuantitativa que describe a toda la población: Media $\mu$, Desviación Estándar $\sigma$, Proporción $\pi$. Son valores fijos pero generalmente desconocidos.
    \item \textbf{Estadístico o Estimador (Letras Latinas):} Medida cuantitativa calculada a partir de los datos de la muestra: Media muestral $\bar{X}$, Desviación estándar muestral $S$, Proporción muestral $p$. Son variables aleatorias cuyo valor varía de una muestra a otra.
\end{itemize}
\end{definicion}

\section{Distribución Muestral de la Media ($\bar{X}$)}
Si a partir de una población se extraen todas las muestras posibles de tamaño $n$, cada muestra generará un promedio $\bar{X}$ particular. La distribución de probabilidad formada por todos estos promedios posibles se denomina \textbf{Distribución Muestral de la Media}.

\subsection{Propiedades de la Distribución Muestral}
\begin{enumerate}
    \item \textbf{Insesgadez de la Media Muestral:}
    El valor esperado de todos los promedios muestrales posibles es exactamente igual a la media poblacional $\mu$:
    \begin{equation}
        E(\bar{X}) = \mu_{\bar{X}} = \mu
    \end{equation}
    
    \item \textbf{Error Estándar de la Media ($\sigma_{\bar{X}}$):}
    Mide la dispersión o variabilidad de los promedios muestrales respecto a la media verdadera:
    \begin{itemize}
        \item \textbf{Población Infinita o Muestreo con Reemplazo:}
        \begin{equation}
            \sigma_{\bar{X}} = \frac{\sigma}{\sqrt{n}}
        \end{equation}
        \item \textbf{Población Finita Sin Reemplazo (cuando $n/N > 0.05$):}
        \begin{equation}
            \sigma_{\bar{X}} = \frac{\sigma}{\sqrt{n}} \sqrt{\frac{N - n}{N - 1}}
        \end{equation}
        donde $\sqrt{\frac{N-n}{N-1}}$ es el \textbf{Factor de Corrección por Población Finita (FCPF)}.
    \end{itemize}
\end{enumerate}

\section{El Teorema del Límite Central (TLC)}
\begin{definicion}{Teorema del Límite Central}
Al seleccionar muestras aleatorias simples de tamaño $n$ de cualquier población (sin importar si la distribución de origen es simétrica, asimétrica, uniforme o desconocida), a medida que el tamaño de la muestra $n$ aumenta, la distribución muestral de la media $\bar{X}$ se aproxima a una \textbf{Distribución Normal} con media $\mu$ y varianza $\sigma^2 / n$:
\begin{equation}
    \bar{X} \xrightarrow{d} N\left( \mu, \, \frac{\sigma^2}{n} \right)
\end{equation}
Para propósitos prácticos en administración y contaduría, la aproximación normal es excelente si:
\begin{equation}
    n \ge 30
\end{equation}
\end{definicion}

\section{Intervalos de Confianza para la Media Poblacional ($\mu$)}
Una \textbf{estimación puntual} proporciona un único valor numérico para estimar $\mu$, pero no indica el grado de error o precisión. Un \textbf{Intervalo de Confianza} establece un rango numérico $[\text{Límite Inferior}, \text{Límite Superior}]$ dentro del cual se espera encontrar el verdadero parámetro poblacional con un nivel de confianza estipulado $(1 - \alpha)$, donde $\alpha$ es el nivel de significancia o riesgo de error.

\subsection{Niveles de Confianza Comunes y Valores Críticos $Z_{\alpha/2}$}
\begin{center}
\begin{tabular}{cccc}
\toprule
\textbf{Nivel de Confianza ($1 - \alpha$)} & \textbf{Significancia ($\alpha$)} & \textbf{Área en dos colas ($\alpha/2$)} & \textbf{Valor Crítico $Z_{\alpha/2}$} \\
\midrule
$90\%$ & $0.10$ & $0.050$ & $1.645$ \\
$95\%$ & $0.05$ & $0.025$ & $1.960$ \\
$99\%$ & $0.01$ & $0.005$ & $2.576$ \\
\bottomrule
\end{tabular}
\end{center}

\subsection{Caso 1: Varianza Poblacional Conocida ($\sigma$) o Muestra Grande ($n \ge 30$)}
La fórmula general del intervalo de confianza al $(1 - \alpha)\%$ es:
\begin{equation}
    \bar{X} \pm Z_{\alpha/2} \cdot \sigma_{\bar{X}} \iff \left[ \bar{X} - Z_{\alpha/2} \frac{\sigma}{\sqrt{n}}, \quad \bar{X} + Z_{\alpha/2} \frac{\sigma}{\sqrt{n}} \right]
\end{equation}
Si $\sigma$ es desconocida pero la muestra es grande ($n \ge 30$), por el TLC se sustituye $\sigma$ por la desviación estándar muestral $S$:
\begin{equation}
    \bar{X} \pm Z_{\alpha/2} \frac{S}{\sqrt{n}}
\end{equation}

\begin{ejemplo}{Auditoría de Saldos de Cuentas por Cobrar (Basado en Anderson)}
Una firma auditora en Santa Cruz revisa los saldos deudores de una empresa comercializadora. A partir de una muestra aleatoria de $n = 64$ cuentas por cobrar, se obtuvo un saldo promedio muestral de $\bar{X} = 3{,}850$ Bs con una desviación estándar muestral de $S = 480$ Bs.
\begin{enumerate}[label=\alph*)]
    \item Construya un intervalo de confianza del $95\%$ para el saldo promedio verdadero de todas las cuentas por cobrar de la empresa.
    \item Construya un intervalo de confianza del $99\%$ para la media poblacional.
    \item Interprete los resultados desde la perspectiva de auditoría financiera.
\end{enumerate}

\tcbline
\textbf{Solución Paso a Paso:}
Datos: $n = 64 \ge 30$, $\bar{X} = 3{,}850$, $S = 480$.
\begin{enumerate}[label=\alph*)]
    \item \textbf{Intervalo al $95\%$ de confianza ($Z_{0.025} = 1.96$):}
    Calculamos el margen de error $E$:
    \[
    E = Z_{\alpha/2} \frac{S}{\sqrt{n}} = (1.96) \frac{480}{\sqrt{64}} = 1.96 \cdot \frac{480}{8} = 1.96 \cdot 60 = 117.60 \text{ Bs.}
    \]
    Construimos los límites:
    \[
    \text{IC}_{95\%} = 3{,}850 \pm 117.60 \implies [3{,}732.40 \text{ Bs}, \quad 3{,}967.60 \text{ Bs}]
    \]
    
    \item \textbf{Intervalo al $99\%$ de confianza ($Z_{0.005} = 2.576$):}
    \[
    E = (2.576)(60) = 154.56 \text{ Bs.}
    \]
    \[
    \text{IC}_{99\%} = 3{,}850 \pm 154.56 \implies [3{,}695.44 \text{ Bs}, \quad 4{,}004.56 \text{ Bs}]
    \]
    
    \item \textbf{Interpretación Contable:} Con un 95\% de confianza estadística, el saldo promedio real adeudado por los clientes de la empresa se encuentra entre \textbf{3,732.40 Bs} y \textbf{3,967.60 Bs}. Si la administración reportó en su balance un saldo medio de 3,800 Bs, la auditoría valida que dicha cifra es plenamente razonable y no presenta evidencia de sobreestimación.
\end{enumerate}
\end{ejemplo}

\subsection{Caso 2: Varianza Desconocida ($\sigma$) y Muestra Pequeña ($n < 30$)}
Cuando la población de origen es normal, la desviación estándar poblacional $\sigma$ es desconocida y la muestra es pequeña ($n < 30$), la variable estandarizada no sigue una normal sino una \textbf{Distribución $t$ de Student} con $\nu = n - 1$ grados de libertad:
\begin{equation}
    \bar{X} \pm t_{\alpha/2, \, n-1} \frac{S}{\sqrt{n}}
\end{equation}

\begin{ejemplo}{Auditoría de Costos de Transporte con Muestra Pequeña (Basado en Lind)}
El auditor de costos de una empresa logística toma una muestra aleatoria de $n = 16$ viajes de transporte de carga pesada hacia la Chiquitanía, obteniendo un costo de combustible promedio de $\bar{X} = 1{,}420$ Bs con $S = 180$ Bs. Suponiendo que los costos tienen distribución normal, construya un intervalo de confianza del $95\%$ para el costo medio poblacional por viaje.

\tcbline
\textbf{Solución:}
\begin{enumerate}
    \item \textbf{Grados de libertad:} $\nu = n - 1 = 16 - 1 = 15$.
    \item \textbf{Valor crítico $t$:} En la tabla $t$ para $\alpha/2 = 0.025$ y $15$ gl: $t_{0.025, 15} = 2.131$.
    \item \textbf{Margen de error $E$:}
    \[
    E = t_{\alpha/2, n-1} \frac{S}{\sqrt{n}} = (2.131) \frac{180}{\sqrt{16}} = 2.131 \cdot \frac{180}{4} = 2.131 \cdot 45 = 95.90 \text{ Bs.}
    \]
    \item \textbf{Intervalo de Confianza:}
    \[
    \text{IC}_{95\%} = 1{,}420 \pm 95.90 \implies [1{,}324.10 \text{ Bs}, \quad 1{,}515.90 \text{ Bs}]
    \]
\end{enumerate}
\end{ejemplo}

\section{Determinación del Tamaño de Muestra Necesario ($n$)}
Uno de los problemas más críticos en la planificación de una auditoría es decidir cuántos elementos deben seleccionarse para garantizar que el margen de error muestral no exceda un límite tolerable $E$ a un nivel de confianza fijado $(1-\alpha)$.

\subsection{Fórmula para Población Infinita o Muy Grande}
Despejando $n$ de la ecuación del margen de error $E = Z_{\alpha/2} \frac{\sigma}{\sqrt{n}}$:
\begin{equation}
    n = \left( \frac{Z_{\alpha/2} \cdot \sigma}{E} \right)^2 = \frac{Z_{\alpha/2}^2 \cdot \sigma^2}{E^2}
\end{equation}

\subsection{Fórmula para Población Finita ($N$ Conocida)}
Cuando se conoce el tamaño total de la población $N$:
\begin{equation}
    n = \frac{N \cdot Z_{\alpha/2}^2 \cdot \sigma^2}{(N - 1) E^2 + Z_{\alpha/2}^2 \cdot \sigma^2}
\end{equation}

\begin{ejemplo}{Cálculo de Muestra para Auditoría Fiscal de Inventarios (Basado en Anderson)}
El socio director de una firma auditora debe planificar la revisión de inventarios de una ferretería industrial que almacena $N = 3{,}500$ ítems distintos. Estudios previos indican que la desviación estándar del valor por ítem es de $\sigma = 250$ Bs. La firma desea estimar el valor medio del inventario con un margen de error máximo admisible de $E = \pm 20$ Bs y un nivel de confianza del $95\%$ ($Z = 1.96$). ¿Cuántos ítems deben auditarse?

\tcbline
\textbf{Solución:}
Datos: $N = 3{,}500$, $\sigma = 250$, $E = 20$, $Z = 1.96$.
\begin{align*}
    n &= \frac{3{,}500 \cdot (1.96)^2 \cdot (250)^2}{(3{,}500 - 1)(20)^2 + (1.96)^2 (250)^2} \\
    &= \frac{3{,}500 \cdot 3.8416 \cdot 62{,}500}{3{,}499 \cdot 400 + 3.8416 \cdot 62{,}500} \\
    &= \frac{840{,}350{,}000}{1{,}399{,}600 + 240{,}100} = \frac{840{,}350{,}000}{1{,}639{,}700} \approx 512.50 \implies \mathbf{513 \text{ ítems.}}
\end{align*}
\textbf{Regla Práctica:} En el cálculo de tamaño de muestra, el resultado fraccionario siempre se redondea al entero superior inmediato para garantizar el nivel de confianza exigido.
\end{ejemplo}

\section{Guía y Resumen de Decisiones Metodológicas}
\begin{center}
\begin{tabularx}{\textwidth}{lcccc}
\toprule
\textbf{Distribución Poblacional} & \textbf{Tamaño Muestra} & \textbf{Varianza $\sigma^2$} & \textbf{Estadístico} & \textbf{Fórmula del Intervalo} \\
\midrule
Cualquiera & Grande ($n \ge 30$) & Conocida & $Z$ & $\bar{X} \pm Z_{\alpha/2} \frac{\sigma}{\sqrt{n}}$ \\
Cualquiera & Grande ($n \ge 30$) & Desconocida & $Z$ & $\bar{X} \pm Z_{\alpha/2} \frac{S}{\sqrt{n}}$ \\
Normal & Pequeña ($n < 30$) & Conocida & $Z$ & $\bar{X} \pm Z_{\alpha/2} \frac{\sigma}{\sqrt{n}}$ \\
Normal & Pequeña ($n < 30$) & Desconocida & $t$ & $\bar{X} \pm t_{\alpha/2, n-1} \frac{S}{\sqrt{n}}$ \\
No Normal & Pequeña ($n < 30$) & Desconocida & -- & Métodos No Paramétricos \\
\bottomrule
\end{tabularx}
\end{center}

\section{Formulario Resumen del Capítulo 5}
\begin{formulario}[Resumen de Fórmulas — Muestreo e Intervalos de Confianza]
\begin{itemize}
    \item \textbf{Error Estándar:} $\sigma_{\bar{X}} = \dfrac{\sigma}{\sqrt{n}}$ \quad (con FCPF: $\sigma_{\bar{X}} = \dfrac{\sigma}{\sqrt{n}} \sqrt{\dfrac{N-n}{N-1}}$)
    \item \textbf{IC con Normal ($Z$):} $\bar{X} \pm Z_{\alpha/2} \dfrac{S}{\sqrt{n}}$
    \item \textbf{IC con $t$ de Student:} $\bar{X} \pm t_{\alpha/2, n-1} \dfrac{S}{\sqrt{n}}$
    \item \textbf{Tamaño de Muestra (Infinita):} $n = \left( \dfrac{Z_{\alpha/2} \cdot \sigma}{E} \right)^2$
    \item \textbf{Tamaño de Muestra (Finita):} $n = \dfrac{N \cdot Z_{\alpha/2}^2 \cdot \sigma^2}{(N-1)E^2 + Z_{\alpha/2}^2 \cdot \sigma^2}$
\end{itemize}
\end{formulario}

\section{Ejercicios Propuestos de la Unidad}

\subsection*{Nivel 1: Conceptos Básicos y Distribución Muestral}
\begin{enumerate}
    \item Explique la importancia del Teorema del Límite Central en la inferencia estadística aplicada a la auditoría y cuándo no es necesario que la población original sea normal.
    \item Una población de cuentas por cobrar tiene una media de $\mu = 4{,}200$ Bs y una desviación estándar $\sigma = 800$ Bs. Si se extrae una muestra aleatoria simple de $n = 64$ cuentas:
    \begin{enumerate}[label=\alph*)]
        \item Calcule el error estándar de la media muestral $\sigma_{\bar{X}}$.
        \item ¿Cuál es la probabilidad de que la media muestral esté entre $4{,}000$ Bs y $4{,}400$ Bs?
        \item ¿Cuál es la probabilidad de que la media muestral supere los $4{,}350$ Bs?
    \end{enumerate}
\end{enumerate}

\subsection*{Nivel 2: Intervalos de Confianza para la Media}
\begin{enumerate}[resume]
    \item Una muestra aleatoria de $n = 50$ facturas comerciales en una empresa importadora arrojó un importe medio de $\bar{X} = 820$ Bs con una desviación estándar muestral $S = 140$ Bs.
    \begin{enumerate}[label=\alph*)]
        \item Construya e interprete un intervalo de confianza del 95\% para el importe medio de todas las facturas de la empresa.
        \item Construya un intervalo de confianza del 99\% para la misma media. Compare la amplitud de ambos intervalos y explique la relación entre nivel de confianza y precisión.
    \end{enumerate}
    \item Un auditor fiscal desea estimar el tiempo medio que toma la liquidación de trámites aduaneros en la frontera. Selecciona una muestra pequeña de $n = 16$ trámites y obtiene una media muestral $\bar{X} = 32.5$ horas con una desviación estándar muestral $S = 6.4$ horas. Asumiendo que los tiempos se distribuyen normalmente:
    \begin{enumerate}[label=\alph*)]
        \item Identifique la distribución apropiada y los grados de libertad.
        \item Calcule el intervalo de confianza del 95\% para el tiempo medio poblacional.
    \end{enumerate}
\end{enumerate}

\subsection*{Nivel 3: Determinación del Tamaño Óptimo de Muestra}
\begin{enumerate}[resume]
    \item La gerencia de operaciones de una entidad bancaria desea estimar el tiempo medio de permanencia de los clientes en plataforma de atención con un margen de error máximo admisible de $E = 1.5$ minutos y un nivel de confianza del 99\% ($Z = 2.576$). Si estudios preliminares indican una desviación estándar de $\sigma = 6$ minutos:
    \begin{enumerate}[label=\alph*)]
        \item ¿Cuántos clientes deben incluirse en la muestra si la población se considera infinita?
        \item Si la sucursal atiende a un total conocido de $N = 800$ clientes por día, ¿cuál será el tamaño de muestra ajustado por población finita?
    \end{enumerate}
    \item En una auditoría de existencias sobre un inventario de $N = 2{,}500$ repuestos automotrices, se requiere estimar el valor promedio por ítem con un error de estimación no mayor a $\pm 15$ Bs con un 95\% de confianza. La desviación estándar estimada del valor unitario es $\sigma = 120$ Bs. Determine el tamaño de muestra requerido aplicando el Factor de Corrección para Poblaciones Finitas (FCPF).
\end{enumerate}


% Capítulo 6: Pruebas de Hipótesis
\chapter{Pruebas de Hipótesis para una Población}
\label{cap:pruebas_hipotesis}

\section*{Objetivos de Aprendizaje de la Unidad}
Al finalizar el estudio de esta unidad, el estudiante de Contaduría Pública será capaz de:
\begin{enumerate}
    \item Formular con precisión la \textbf{hipótesis nula ($H_0$)} y la \textbf{hipótesis alternativa ($H_1$)} para problemas de decisión en auditoría, finanzas y control de gestión.
    \item Comprender y evaluar el riesgo de cometer un \textbf{Error Tipo I ($\alpha$)} y un \textbf{Error Tipo II ($\beta$)}, así como el concepto de nivel de significancia.
    \item Ejecutar la metodología sistemática de \textbf{5 pasos} para pruebas de hipótesis bilaterales y unilaterales (cola izquierda y cola derecha).
    \item Realizar pruebas de hipótesis para la \textbf{media poblacional ($\mu$)} utilizando la distribución Normal ($Z$) y la distribución $t$ de Student.
    \item Tomar decisiones mediante el enfoque tradicional de \textbf{valores críticos} y el enfoque contemporáneo del \textbf{valor $p$ ($p$-value)}.
    \item Efectuar pruebas de hipótesis para la \textbf{proporción poblacional ($\pi$)} en auditoría tributaria y cumplimiento normativo.
\end{enumerate}

\section{Fundamentos de la Prueba de Hipótesis en la Toma de Decisiones}
En la práctica profesional del contador público y del auditor financiero, con frecuencia es indispensable tomar decisiones sobre afirmaciones relativas a parámetros poblacionales basadas únicamente en la evidencia proporcionada por una muestra.

Por ejemplo, la gerencia financiera puede afirmar que el tiempo medio de cobranza de facturas es de 30 días, o el departamento tributario puede sostener que la tasa de error en las declaraciones juradas no supera el 3\%. La \textbf{Prueba de Hipótesis} es el procedimiento formal basado en la teoría de la probabilidad que permite determinar si la evidencia muestral respalda o refuta dichas aseveraciones.

\begin{definicion}{Hipótesis Estadística, Hipótesis Nula e Hipótesis Alternativa}
\begin{itemize}
    \item \textbf{Hipótesis Estadística:} Afirmación o conjetura acerca del valor numérico de uno o más parámetros de una población.
    \item \textbf{Hipótesis Nula ($H_0$):} Enunciado tentativo que establece que el parámetro es igual a un valor específico o que no existe diferencia ni efecto real. Siempre contiene el signo de igualdad ($=$, $\le$ o $\ge$). Es la hipótesis que se somete a prueba directa y que se mantiene a menos que los datos muestrales aporten evidencia contundente en su contra.
    \item \textbf{Hipótesis Alternativa ($H_1$ o $H_a$):} Enunciado que se acepta como verdadero si la evidencia muestral conduce al rechazo de $H_0$. Contiene signos de desigualdad estricta ($\ne$, $<$, $>$). Suele representar la sospecha del auditor o la nueva hipótesis de investigación.
\end{itemize}
\end{definicion}

\section{Estructura de las Pruebas de Hipótesis y Tipos de Errores}

\subsection{Tipos de Pruebas según la Dirección de la Hipótesis Alternativa}
\begin{enumerate}
    \item \textbf{Prueba Bilateral o de Dos Colas:}
    \[
    \begin{cases} H_0: \mu = \mu_0 \\ H_1: \mu \ne \mu_0 \end{cases}
    \]
    La región de rechazo se divide en partes iguales ($\alpha/2$) en ambos extremos de la distribución.
    
    \item \textbf{Prueba Unilateral de Cola Derecha (Cola Superior):}
    \[
    \begin{cases} H_0: \mu \le \mu_0 \\ H_1: \mu > \mu_0 \end{cases}
    \]
    La región de rechazo ($\alpha$) se concentra totalmente en el extremo derecho superior.
    
    \item \textbf{Prueba Unilateral de Cola Izquierda (Cola Inferior):}
    \[
    \begin{cases} H_0: \mu \ge \mu_0 \\ H_1: \mu < \mu_0 \end{cases}
    \]
    La región de rechazo ($\alpha$) se concentra totalmente en el extremo izquierdo inferior.
\end{enumerate}

\subsection{Matriz de Decisiones y Errores Tipo I y Tipo II}
Al contrastar una hipótesis estadística con base en una muestra aleatoria, existe la posibilidad de tomar una decisión incorrecta debido a la variabilidad del muestreo:

\begin{center}
\small
\begin{tabular}{l|c|c}
\toprule
\multirow{2}{*}{\textbf{Decisión Muestral}} & \multicolumn{2}{c}{\textbf{Estado Real de la Población}} \\
\cmidrule{2-3}
& \textbf{$H_0$ es Verdadera} & \textbf{$H_0$ es Falsa} \\
\midrule
\textbf{No Rechazar $H_0$} & Decisión Correcta ($1 - \alpha$) & Error Tipo II ($\beta$) \\
\textbf{Rechazar $H_0$} & \textbf{Error Tipo I ($\alpha$)} & Decisión Correcta: Potencia ($1 - \beta$) \\
\bottomrule
\end{tabular}
\end{center}

\begin{definicion}{Error Tipo I, Nivel de Significancia y Error Tipo II}
\begin{itemize}
    \item \textbf{Error Tipo I ($\alpha$):} Rechazar la hipótesis nula $H_0$ cuando en realidad es verdadera. La probabilidad de cometer este error se denomina \textbf{Nivel de Significancia ($\alpha$)}, y típicamente se fija en $0.05$ (5\%), $0.01$ (1\%) o $0.10$ (10\%).
    \item \textbf{Error Tipo II ($\beta$):} No rechazar la hipótesis nula $H_0$ cuando en realidad es falsa.
    \item \textbf{Potencia de la Prueba ($1 - \beta$):} Probabilidad de rechazar correctamente una hipótesis nula que es falsa.
\end{itemize}
\end{definicion}

\begin{alerta}[Riesgo del Auditor y Errores en Auditoría]
En una auditoría financiera:
\begin{itemize}
    \item Un \textbf{Error Tipo I} significa concluir que los estados financieros contienen errores materiales cuando en realidad son correctos (genera costos adicionales innecesarios de revisión).
    \item Un \textbf{Error Tipo II} (Riesgo de Detección) significa no detectar incorrecciones materiales y emitir una opinión favorable indebida, lo cual acarrea seria responsabilidad legal para el auditor.
\end{itemize}
\end{alerta}

\section{Metodología General de 5 Pasos para la Prueba de Hipótesis}
Siguiendo las pautas pedagógicas de \textit{Anderson, Sweeney \& Williams} y \textit{Lind, Marchal \& Wathen}, todo contraste de hipótesis debe desarrollarse conforme a los siguientes cinco pasos:

\begin{formulario}[Los 5 Pasos Metodológicos de una Prueba de Hipótesis]
\begin{enumerate}
    \item \textbf{Paso 1:} Formular la Hipótesis Nula ($H_0$) y la Hipótesis Alternativa ($H_1$).
    \item \textbf{Paso 2:} Seleccionar el Nivel de Significancia ($\alpha$).
    \item \textbf{Paso 3:} Identificar el Estadístico de Prueba ($Z$ o $t$) y su distribución muestral.
    \item \textbf{Paso 4:} Formular la Regla de Decisión (mediante Valor Crítico o Enfoque del Valor $p$).
    \item \textbf{Paso 5:} Calcular el estadístico muestral, tomar la decisión y redactar la conclusión en el contexto contable/financiero.
\end{enumerate}
\end{formulario}

\section{Pruebas para la Media Poblacional ($\mu$)}

\subsection{Caso A: Desviación Estándar Poblacional $\sigma$ Conocida (Distribución $Z$)}
Cuando la desviación estándar poblacional $\sigma$ es conocida (o la muestra es grande $n \ge 30$ por el TLC), el estadístico de prueba es:
\begin{equation}
    Z_{\text{calc}} = \frac{\bar{X} - \mu_0}{\frac{\sigma}{\sqrt{n}}}
\end{equation}

\begin{ejemplo}{Verificación del Saldo Promedio en Cuentas por Cobrar (Ref. Anderson)}
El departamento de crédito y cobranzas de una distribuidora comercial en Santa Cruz afirma que el saldo promedio adeudado por sus clientes es de $3{,}500$ Bs con una desviación estándar histórica $\sigma = 600$ Bs. 

Un auditor externo selecciona una muestra aleatoria de $n = 64$ cuentas y encuentra una media muestral de $\bar{X} = 3{,}320$ Bs. Con un nivel de significancia de $\alpha = 0.05$, ¿existe evidencia suficiente para refutar la afirmación de la gerencia?

\tcbline
\textbf{Solución Paso a Paso:}
\begin{enumerate}
    \item \textbf{Paso 1: Planteamiento de Hipótesis (Prueba Bilateral):}
    \[
    \begin{cases} H_0: \mu = 3{,}500 \text{ Bs} \\ H_1: \mu \ne 3{,}500 \text{ Bs} \end{cases}
    \]
    
    \item \textbf{Paso 2: Nivel de Significancia:} $\alpha = 0.05 \implies \alpha/2 = 0.025$.
    
    \item \textbf{Paso 3: Estadístico de Prueba:}
    Como $\sigma = 600$ es conocida y $n = 64 \ge 30$:
    \[
    Z_{\text{calc}} = \frac{\bar{X} - \mu_0}{\frac{\sigma}{\sqrt{n}}}
    \]
    
    \item \textbf{Paso 4: Regla de Decisión (Valores Críticos):}
    Para $\alpha/2 = 0.025$, los valores críticos en la tabla Normal Estándar son $Z_{\text{crit}} = \pm 1.96$.
    \begin{itemize}
        \item Se rechaza $H_0$ si $Z_{\text{calc}} < -1.96$ o si $Z_{\text{calc}} > +1.96$.
        \item No se rechaza $H_0$ si $-1.96 \le Z_{\text{calc}} \le +1.96$.
    \end{itemize}
    
    \item \textbf{Paso 5: Cálculo del Estadístico y Decisión:}
    \[
    \sigma_{\bar{X}} = \frac{600}{\sqrt{64}} = \frac{600}{8} = 75 \text{ Bs}
    \]
    \[
    Z_{\text{calc}} = \frac{3{,}320 - 3{,}500}{75} = \frac{-180}{75} = -2.40
    \]
    Como $Z_{\text{calc}} = -2.40 < -1.96$, el estadístico cae en la \textbf{región de rechazo}.
    
    \textbf{Enfoque del Valor $p$ ($p$-value):}
    \[
    p\text{-valor} = 2 \cdot P(Z \le -2.40) = 2 \cdot (0.0082) = 0.0164
    \]
    Como $p\text{-valor} = 0.0164 < \alpha = 0.05$, se confirma el rechazo de $H_0$.
\end{enumerate}

\textbf{Conclusión Contable:} Con un 5\% de nivel de significancia (95\% de confianza), existe evidencia estadística suficiente para rechazar la afirmación de la empresa. El saldo promedio real de las cuentas por cobrar es significativamente diferente de $3{,}500$ Bs (de hecho, es inferior), lo cual debe reflejarse en los informes de auditoría.
\end{ejemplo}

\subsection{Caso B: Desviación Estándar Poblacional $\sigma$ Desconocida y Muestra Pequeña (Distribución $t$ de Student)}
Cuando la población es normal, $\sigma$ es desconocida y la muestra es pequeña ($n < 30$), se utiliza la desviación estándar muestral $S$ y la \textbf{Distribución $t$ de Student} con $gl = n - 1$ grados de libertad:
\begin{equation}
    t_{\text{calc}} = \frac{\bar{X} - \mu_0}{\frac{S}{\sqrt{n}}}, \quad gl = n - 1
\end{equation}

\begin{casocontable}{Auditoría de Tiempos de Cierre Mensual Contable (Ref. Lind)}
La firma de auditoría «Contadores & Asociados» implementó un nuevo software ERP con la promesa de que el tiempo promedio necesario para procesar el cierre contable mensual sería \textbf{menor a 16 horas}. 

Para comprobarlo, se auditó una muestra de $n = 16$ cierres contables mensuales en diversas empresas clientes, obteniéndose una media muestral de $\bar{X} = 14.5$ horas y una desviación estándar muestral $S = 2.8$ horas. Con $\alpha = 0.01$, ¿el nuevo sistema cumple su promesa?

\tcbline
\textbf{Solución Paso a Paso:}
\begin{enumerate}
    \item \textbf{Paso 1: Planteamiento de Hipótesis (Cola Izquierda):}
    \[
    \begin{cases} H_0: \mu \ge 16 \text{ horas} \\ H_1: \mu < 16 \text{ horas} \text{ (el ERP reduce el tiempo)} \end{cases}
    \]
    
    \item \textbf{Paso 2: Nivel de Significancia:} $\alpha = 0.01$.
    
    \item \textbf{Paso 3: Estadístico de Prueba:}
    Distribución $t$ de Student con $gl = n - 1 = 16 - 1 = 15$ grados de libertad.
    
    \item \textbf{Paso 4: Regla de Decisión (Valor Crítico):}
    Buscando en la tabla $t$ con $gl = 15$ y $\alpha = 0.01$ (unilateral izquierda):
    \[
    t_{\text{crit}} = -t_{0.01, 15} = -2.602
    \]
    Se rechaza $H_0$ si $t_{\text{calc}} < -2.602$.
    
    \item \textbf{Paso 5: Cálculo del Estadístico y Decisión:}
    \[
    S_{\bar{X}} = \frac{S}{\sqrt{n}} = \frac{2.8}{\sqrt{16}} = \frac{2.8}{4} = 0.70 \text{ horas}
    \]
    \[
    t_{\text{calc}} = \frac{14.5 - 16.0}{0.70} = \frac{-1.50}{0.70} \approx -2.143
    \]
    Como $t_{\text{calc}} = -2.143 > -2.602$, el estadístico \textbf{NO cae en la región de rechazo}.
\end{enumerate}

\textbf{Conclusión para la Dirección:} Con un nivel de significancia de $\alpha = 0.01$, no existe evidencia estadística suficiente para rechazar $H_0$. Aunque el promedio muestral fue de 14.5 horas, la variabilidad de los datos no permite afirmar con 99\% de certeza que el tiempo medio poblacional sea estrictamente inferior a 16 horas.
\end{casocontable}

\section{Pruebas de Hipótesis para la Proporción Poblacional ($\pi$)}
En auditoría tributaria y evaluación de cumplimiento, la variable de interés frecuentemente es un atributo dicotómico (ej. factura con/sin retención del IUE/IT, cuenta al día/en mora). 

Para muestras donde $n \cdot \pi_0 \ge 5$ y $n \cdot (1 - \pi_0) \ge 5$, la distribución muestral de la proporción $p$ se aproxima a la normal:
\begin{equation}
    Z_{\text{calc}} = \frac{p - \pi_0}{\sigma_p} = \frac{p - \pi_0}{\sqrt{\frac{\pi_0 (1 - \pi_0)}{n}}}
\end{equation}

\begin{ejemplo}{Auditoría Tributaria de Emisión de Facturas Electrónicas}
La administración tributaria nacional estipula que una empresa no debe superar una tasa de error del $\pi_0 = 4\%$ ($0.04$) en la emisión de facturas electrónicas para mantener su certificación de contribuyente preferente.

En una auditoría impositiva se revisa una muestra aleatoria de $n = 400$ facturas emitidas por la compañía durante el último trimestre, detectándose $x = 24$ facturas con errores de codificación o timbrado. Con un nivel de significancia de $\alpha = 0.05$, ¿se debe revocar la certificación a la empresa?

\tcbline
\textbf{Solución Paso a Paso:}
\begin{enumerate}
    \item \textbf{Paso 1: Hipótesis (Prueba Unilateral de Cola Derecha):}
    \[
    \begin{cases} H_0: \pi \le 0.04 \text{ (cumple el estándar)} \\ H_1: \pi > 0.04 \text{ (supera la tasa de error permitida)} \end{cases}
    \]
    
    \item \textbf{Paso 2: Nivel de Significancia:} $\alpha = 0.05$.
    
    \item \textbf{Paso 3: Proporción Muestral y Estadístico:}
    \[
    p = \frac{x}{n} = \frac{24}{400} = 0.06 \quad (6\%)
    \]
    Verificación de supuestos: $400(0.04) = 16 \ge 5$ y $400(0.96) = 384 \ge 5$.
    
    \item \textbf{Paso 4: Regla de Decisión:}
    Para $\alpha = 0.05$ (cola derecha): $Z_{\text{crit}} = +1.645$.
    Se rechaza $H_0$ si $Z_{\text{calc}} > +1.645$.
    
    \item \textbf{Paso 5: Cálculo y Decisión:}
    \[
    \sigma_p = \sqrt{\frac{0.04(1 - 0.04)}{400}} = \sqrt{\frac{0.04 \cdot 0.96}{400}} = \sqrt{\frac{0.0384}{400}} = \sqrt{0.000096} \approx 0.009798
    \]
    \[
    Z_{\text{calc}} = \frac{0.06 - 0.04}{0.009798} = \frac{0.02}{0.009798} \approx +2.04
    \]
    Como $Z_{\text{calc}} = +2.04 > +1.645$, se \textbf{rechaza $H_0$}.
    
    \[
    p\text{-valor} = P(Z \ge 2.04) = 1 - 0.9793 = 0.0207 < 0.05
    \]
\end{enumerate}

\textbf{Dictamen del Auditor:} A un nivel de significancia del 5\%, la evidencia muestral demuestra que la proporción real de facturas con errores supera el 4\% máximo permitido. Se recomienda a la gerencia adoptar acciones correctivas de inmediato en el sistema de facturación.
\end{ejemplo}

\section{Ejercicios Propuestos de la Unidad}

\subsection*{Nivel 1: Conceptos y Mecánica Operativa}
\begin{enumerate}
    \item Enuncie las diferencias conceptuales entre un Error Tipo I y un Error Tipo II en pruebas estadísticas. Dé un ejemplo de cada uno en el contexto del peritaje contable.
    \item Para cada una de las siguientes situaciones, formule la hipótesis nula ($H_0$) y la hipótesis alternativa ($H_1$):
    \begin{enumerate}
        \item Un banco afirma que el saldo medio de sus depósitos a plazo fijo es de $50{,}000$ Bs.
        \item Un auditor sospecha que el tiempo medio de atención a clientes en caja supera los 8 minutos.
        \item Un gerente de compras sostiene que menos del 2\% de los insumos recibidos presentan defectos.
    \end{enumerate}
    \item Si se realiza una prueba de hipótesis bilateral para la media con muestra grande y se obtiene un estadístico $Z_{\text{calc}} = -2.15$, determine la decisión que se debe tomar con niveles de significancia de:
    \begin{enumerate}
        \item $\alpha = 0.05$
        \item $\alpha = 0.01$
    \end{enumerate}
\end{enumerate}

\subsection*{Nivel 2: Aplicaciones en Finanzas y Auditoría}
\begin{enumerate}[resume]
    \item Una cadena de farmacias en Santa Cruz sostiene que el gasto promedio por cliente en cada compra es de $85$ Bs con una desviación estándar $\sigma = 18$ Bs. Un analista financiero toma una muestra de $n = 81$ compras y registra una media muestral de $\bar{X} = 81.5$ Bs. Con un nivel de significancia de $\alpha = 0.05$, pruebe si el gasto medio es significativamente menor que $85$ Bs.
    \item Un departamento de auditoría interna de una entidad financiera revisa una muestra de $n = 25$ créditos comerciales para verificar la tasa de interés promedio cobrada. La media muestral fue de $\bar{X} = 14.2\%$ anual con una desviación estándar muestral $S = 1.5\%$. Si la política institucional estipula una tasa media de $15.0\%$, verifique con $\alpha = 0.02$ si la tasa real cobrada difiere de la política oficial.
    \item Una empresa importadora afirma que a lo sumo el 5\% de sus productos importados sufren demoras aduaneras. En una auditoría de comercio exterior se revisan $n = 300$ despachos, encontrándose que 24 tuvieron demoras. Pruebe la afirmación de la empresa con $\alpha = 0.01$.
\end{enumerate}

\subsection*{Nivel 3: Casos Integrales de Toma de Decisiones}
\begin{enumerate}[resume]
    \item La gerencia de recursos humanos de una corporación financiera implementó un programa intensivo de capacitación en Excel avanzado y Power BI para sus analistas contables. Antes del programa, el tiempo medio empleado en elaborar el informe de gestión mensual era de 40 horas con $\sigma = 6$ horas. Tras la capacitación, una muestra aleatoria de $n = 36$ analistas registró una media de $\bar{X} = 37.2$ horas.
    \begin{enumerate}
        \item Formule las hipótesis correspondientes para evaluar si la capacitación redujo efectivamente el tiempo de elaboración.
        \item Calcule el estadístico de prueba y determine el valor $p$ ($p$-value).
        \item ¿Qué recomendación emitiría a la gerencia general con $\alpha = 0.01$?
    \end{enumerate}
    \item En una auditoría forense sobre $n = 500$ transacciones bancarias sospechosas, se identificó que 45 correspondían a transferencias sin respaldo documental completo. Pruebe si la proporción real de transferencias no respaldadas excede el umbral de tolerancia crítica fijado por la Unidad de Investigaciones Financieras (UIF) en $\pi = 0.06$, empleando $\alpha = 0.05$.
\end{enumerate}


% Capítulo 7: Regresión Lineal y Correlación
\chapter{Regresión Lineal Simple y Correlación Aplicada a Costos y Finanzas}
\label{cap:regresion_correlacion}

\section*{Objetivos de Aprendizaje de la Unidad}
Al finalizar el estudio de esta unidad, el estudiante de Contaduría Pública será capaz de:
\begin{enumerate}
    \item Construir e interpretar \textbf{diagramas de dispersión} para identificar patrones de relación lineal, no lineal o ausencia de correlación entre variables empresariales.
    \item Calcular e interpretar con exactitud el \textbf{Coeficiente de Correlación de Pearson ($r$)} y evaluar la fuerza y sentido de la asociación bivariada.
    \item Ajustar la ecuación de la \textbf{recta de regresión lineal simple} mediante el método de los \textbf{Mínimos Cuadrados Ordinarios (MCO)}.
    \item Separar e interpretar los componentes de \textbf{Costo Fijo ($\beta_0$)} y \textbf{Costo Variable Unitario ($\beta_1$)} en la estructura de costos mixtos de una empresa.
    \item Evaluar la bondad de ajuste del modelo mediante el \textbf{Coeficiente de Determinación ($R^2$)} y el \textbf{Error Estándar de la Estimación ($s_{yx}$)}.
    \item Generar \textbf{pronósticos puntuales y por intervalo} para la toma de decisiones gerenciales, financieras y presupuestarias.
\end{enumerate}

\section{Introducción al Análisis de Relaciones Bivariadas}
En las ciencias contables, financieras y administrativas, las decisiones rara vez involucran una sola variable aislada. Por el contrario, los directores financieros y contadores analizan continuamente cómo una variable influye o se relaciona con otra:
\begin{itemize}
    \item ¿Cómo varía el \textbf{Costo Total de Producción ($Y$)} en función del \textbf{Número de Horas-Máquina o Unidades Producidas ($X$)}?
    \item ¿En qué medida el incremento en los \textbf{Gastos de Publicidad ($X$)} impulsa los \textbf{Ingresos por Ventas ($Y$)}?
    \item ¿Cómo influye el \textbf{Plazo de Crédito otorgado ($X$)} en el \textbf{Monto de Cuentas Incobrables ($Y$)}?
\end{itemize}

El \textbf{Análisis de Correlación} cuantifica la fuerza y dirección de la relación lineal entre dos variables, mientras que el \textbf{Análisis de Regresión} formula una ecuación matemática para estimar o predecir el valor de la variable dependiente $Y$ a partir de la variable independiente $X$.

\section{El Diagrama de Dispersión}

\begin{definicion}{Diagrama de Dispersión}
Un \textbf{Diagrama de Dispersión} es una representación gráfica cartesiana donde cada observación muestral se grafica como un punto en el plano bidimensional $(x_i, y_i)$.
\begin{itemize}
    \item \textbf{Eje Horizontal ($X$):} Variable independiente, predictora o explicativa.
    \item \textbf{Eje Vertical ($Y$):} Variable dependiente, de respuesta o pronosticada.
\end{itemize}
\end{definicion}

El diagrama de dispersión es el primer paso indispensable en cualquier estudio econométrico o contable, pues permite verificar visualmente si la relación tiene tendencia lineal positiva, lineal negativa, curvilínea o si los puntos se encuentran completamente dispersos sin correlación aparente.

\section{Coeficiente de Correlación Lineal de Pearson ($r$)}
Desarrollado por Karl Pearson, mide el grado de asociación lineal entre dos variables cuantitativas continuas.

\begin{definicion}{Coeficiente de Correlación Muestral de Pearson ($r$)}
Para una muestra de $n$ pares de observaciones $(x_1, y_1), (x_2, y_2), \dots, (x_n, y_n)$, el coeficiente $r$ se calcula como:
\begin{equation}
    r = \frac{\sum (x_i - \bar{x})(y_i - \bar{y})}{\sqrt{\sum (x_i - \bar{x})^2 \cdot \sum (y_i - \bar{y})^2}} = \frac{n \sum x_i y_i - (\sum x_i)(\sum y_i)}{\sqrt{\left[ n \sum x_i^2 - (\sum x_i)^2 \right] \left[ n \sum y_i^2 - (\sum y_i)^2 \right]}}
\end{equation}
\end{definicion}

\subsection{Propiedades e Interpretación de $r$}
\begin{enumerate}
    \item \textbf{Rango Acotado:} $-1 \le r \le +1$.
    \item \textbf{Sentido de la Relación:}
    \begin{itemize}
        \item Si $r > 0$: Correlación positiva o directa (al aumentar $X$, tiende a aumentar $Y$).
        \item Si $r < 0$: Correlación negativa o inversa (al aumentar $X$, tiende a disminuir $Y$).
        \item Si $r = 0$: Ausencia de correlación lineal (no existe relación lineal entre las variables).
    \end{itemize}
    \item \textbf{Fuerza de la Asociación (Criterio Estándar en Negocios):}
    \begin{itemize}
        \item $|r| \ge 0.90$: Correlación lineal muy fuerte o excelente.
        \item $0.70 \le |r| < 0.90$: Correlación lineal fuerte.
        \item $0.40 \le |r| < 0.70$: Correlación lineal moderada.
        \item $0.20 \le |r| < 0.40$: Correlación lineal débil.
        \item $|r| < 0.20$: Correlación lineal muy débil o nula.
    \end{itemize}
\end{enumerate}

\begin{alerta}[Correlación no Implica Causalidad]
Un valor alto de $r$ demuestra que dos variables numéricas covarían juntas en el tiempo o espacio, pero \textbf{NO prueba que $X$ sea la causa directa de $Y$}. La causalidad solo puede establecerse mediante teoría económica, conocimiento contable del negocio y diseño experimental controlado.
\end{alerta}

\section{Modelo de Regresión Lineal Simple por Mínimos Cuadrados}
El modelo de regresión lineal simple asume que la variable dependiente $Y$ está relacionada linealmente con la variable independiente $X$ según la ecuación poblacional:
\begin{equation}
    Y = \beta_0 + \beta_1 X + \varepsilon
\end{equation}
donde $\beta_0$ es el intercepto poblacional, $\beta_1$ es la pendiente poblacional y $\varepsilon$ es el término de error aleatorio con media cero ($\varepsilon \sim N(0, \sigma^2)$).

A partir de una muestra, se estima la **Ecuación de Regresión Muestral**:
\begin{equation}
    \hat{Y} = b_0 + b_1 X
\end{equation}
donde $\hat{Y}$ es el valor estimado de $Y$, $b_0$ es el estimador del intercepto y $b_1$ es el estimador de la pendiente.

\subsection{Fórmulas de los Mínimos Cuadrados Ordinarios (MCO)}
El método de MCO minimiza la suma de los residuos cuadráticos $\sum (y_i - \hat{y}_i)^2$:

\begin{formulario}[Ecuaciones de Regresión Lineal por MCO]
\begin{equation}
    \mathbf{b_1 = \frac{n \sum x_i y_i - (\sum x_i)(\sum y_i)}{n \sum x_i^2 - (\sum x_i)^2} = \frac{\sum (x_i - \bar{x})(y_i - \bar{y})}{\sum (x_i - \bar{x})^2}}
\end{equation}
\begin{equation}
    \mathbf{b_0 = \bar{y} - b_1 \bar{x} = \frac{\sum y_i - b_1 \sum x_i}{n}}
\end{equation}
\end{formulario}

\subsection{Interpretación Económica y Contable de los Coeficientes}
\begin{itemize}
    \item \textbf{Pendiente ($b_1$):} Representa el cambio promedio estimado en la variable dependiente $Y$ por cada incremento de una unidad en la variable independiente $X$. En contabilidad de costos, $b_1$ representa el **Costo Variable Unitario**.
    \item \textbf{Intercepto ($b_0$):} Representa el valor esperado de $Y$ cuando $X = 0$. En contabilidad de costos, $b_0$ representa el **Costo Fijo Total**.
\end{itemize}

\section{Bondad de Ajuste: Coeficiente de Determinación ($R^2$)}

\begin{definicion}{Coeficiente de Determinación ($R^2$)}
El \textbf{Coeficiente de Determinación ($R^2$)} es la proporción de la variación total de la variable dependiente $Y$ que es explicada por la recta de regresión lineal en función de $X$:
\begin{equation}
    R^2 = (r)^2 = \frac{\text{Variación Explicada (SCR)}}{\text{Variación Total (SCT)}} = 1 - \frac{\text{Variación Residual (SCE)}}{\text{SCT}}
\end{equation}
\end{definicion}

\textbf{Interpretación:} Si $R^2 = 0.88$, se concluye que el \textbf{88\% de la variabilidad observada en $Y$} es explicada por el modelo lineal en función de $X$, mientras que el 12\% restante se debe a factores no considerados o fluctuaciones aleatorias.

\section{Error Estándar de la Estimación ($s_{yx}$)}
Mide la dispersión promedio de los datos observados alrededor de la recta de regresión:
\begin{equation}
    s_{yx} = \sqrt{\frac{\sum (y_i - \hat{y}_i)^2}{n - 2}} = \sqrt{\frac{\sum y_i^2 - b_0 \sum y_i - b_1 \sum x_i y_i}{n - 2}}
\end{equation}

\section{Ejemplos y Casos Resueltos Paso a Paso}

\begin{casocontable}{Separación de Costos Semivariables en una Empresa Industrial (Ref. Anderson)}
El departamento de contabilidad de costos de una empresa manufacturera en el Parque Industrial de Santa Cruz desea desagregar su costo de mantenimiento en componentes fijos y variables. Se recolectaron los datos mensuales de **Horas de Máquina ($X$, en cientos)** y el **Costo Total de Mantenimiento ($Y$, en miles de Bs)** durante los últimos 6 meses:

\begin{center}
\small
\begin{tabular}{c|cc|ccc}
\toprule
\textbf{Mes ($i$)} & \textbf{Horas-Máquina ($x_i$)} & \textbf{Costo Mantenimiento ($y_i$)} & $x_i^2$ & $y_i^2$ & $x_i y_i$ \\
\midrule
1 & 3 & 8 & 9 & 64 & 24 \\
2 & 4 & 9 & 16 & 81 & 36 \\
3 & 6 & 12 & 36 & 144 & 72 \\
4 & 8 & 15 & 64 & 225 & 120 \\
5 & 10 & 17 & 100 & 289 & 170 \\
6 & 11 & 19 & 121 & 361 & 209 \\
\midrule
\textbf{Suma ($\sum$)} & \textbf{42} & \textbf{80} & \textbf{346} & \textbf{1,164} & \textbf{631} \\
\bottomrule
\end{tabular}
\end{center}

Se solicita al auditor:
\begin{enumerate}
    \item Calcular el coeficiente de correlación de Pearson ($r$) e interpretarlo.
    \item Ajustar la recta de regresión de costos por MCO: $\hat{Y} = b_0 + b_1 X$.
    \item Interpretar los coeficientes $b_0$ y $b_1$ en términos de Costo Fijo y Costo Variable.
    \item Calcular el coeficiente de determinación ($R^2$).
    \item Pronosticar el costo total de mantenimiento para un mes con $X = 9$ (900 horas-máquina).
\end{enumerate}

\tcbline
\textbf{Solución Paso a Paso:}
\begin{enumerate}
    \item \textbf{Cálculo de Medias Muestrales:}
    \[
    n = 6, \quad \bar{x} = \frac{\sum x_i}{n} = \frac{42}{6} = 7, \quad \bar{y} = \frac{\sum y_i}{n} = \frac{80}{6} \approx 13.333
    \]
    
    \item \textbf{Cálculo del Coeficiente de Correlación ($r$):}
    \[
    r = \frac{6(631) - (42)(80)}{\sqrt{\left[ 6(346) - (42)^2 \right] \left[ 6(1164) - (80)^2 \right]}}
    \]
    \[
    r = \frac{3786 - 3360}{\sqrt{\left[ 2076 - 1764 \right] \left[ 6984 - 6400 \right]}} = \frac{426}{\sqrt{312 \cdot 584}} = \frac{426}{\sqrt{182208}} = \frac{426}{426.858} \approx 0.998
    \]
    \textbf{Interpretación:} Existe una correlación lineal positiva casi perfecta ($r = +0.998$) entre las horas-máquina y el costo de mantenimiento.
    
    \item \textbf{Cálculo de los Coeficientes de Regresión ($b_1$ y $b_0$):}
    \[
    b_1 = \frac{n \sum xy - (\sum x)(\sum y)}{n \sum x^2 - (\sum x)^2} = \frac{426}{312} \approx 1.3654 \text{ mil Bs por cada 100 horas}
    \]
    \[
    b_0 = \bar{y} - b_1 \bar{x} = 13.3333 - (1.36538)(7) = 13.3333 - 9.5577 = 3.7756 \text{ miles de Bs}
    \]
    
    La ecuación de regresión estimada es:
    \[
    \mathbf{\hat{Y} = 3.776 + 1.365 X}
    \]
    
    \item \textbf{Interpretación Contable de la Estructura de Costos:}
    \begin{itemize}
        \item \textbf{Costo Fijo Base ($b_0$):} $3{,}776$ Bs mensuales independientemente del uso de maquinaria (sueldos fijos de supervisión, seguros y mantenimiento preventivo base).
        \item \textbf{Costo Variable Unitario ($b_1$):} $13.65$ Bs por cada hora-máquina adicional operada ($1.365$ mil Bs por cada 100 horas).
    \end{itemize}
    
    \item \textbf{Coeficiente de Determinación ($R^2$):}
    \[
    R^2 = (r)^2 = (0.998)^2 \approx 0.996 \quad (99.6\%)
    \]
    El \textbf{99.6\% de la variabilidad del costo de mantenimiento} queda explicada por el volumen de horas-máquina operadas. El modelo es extraordinariamente confiable.
    
    \item \textbf{Pronóstico Presupuestario para $X = 9$ (900 horas):}
    \[
    \hat{Y} = 3.776 + 1.365(9) = 3.776 + 12.285 = 16.061 \text{ miles de Bs}
    \]
    \textbf{Conclusión para el Presupuesto:} Para un nivel de actividad de 900 horas-máquina, se debe presupuestar un costo total de mantenimiento de \textbf{$16{,}061$ Bs}.
\end{enumerate}
\end{casocontable}

\begin{ejemplo}{Inversión Publicitaria e Ingresos por Ventas (Ref. Lind)}
Una cadena de tiendas comerciales evaluó el efecto del gasto publicitario semanal ($X$, en miles de Bs) sobre el volumen de ventas ($Y$, en miles de Bs) durante 5 semanas:
$(2, 20), (3, 25), (5, 34), (6, 38), (8, 48)$.

\tcbline
\textbf{Solución Paso a Paso:}
\begin{enumerate}
    \item \textbf{Sumatorias Básicas:}
    $n = 5$, $\sum x = 24$, $\sum y = 165$, $\sum x^2 = 138$, $\sum y^2 = 5909$, $\sum xy = 899$.
    \[
    \bar{x} = \frac{24}{5} = 4.8, \quad \bar{y} = \frac{165}{5} = 33.0
    \]
    
    \item \textbf{Cálculo de Parámetros por MCO:}
    \[
    b_1 = \frac{5(899) - (24)(165)}{5(138) - (24)^2} = \frac{4495 - 3960}{690 - 576} = \frac{535}{114} \approx 4.693
    \]
    \[
    b_0 = 33.0 - 4.693(4.8) = 33.0 - 22.526 = 10.474
    \]
    Ecuación de ventas: $\mathbf{\hat{Y} = 10.474 + 4.693 X}$.
    
    \item \textbf{Interpretación Financiera:}
    Por cada $1{,}000$ Bs adicionales invertidos en publicidad, las ventas aumentan en promedio en $4{,}693$ Bs. Las ventas base sin publicidad son de aproximadamente $10{,}474$ Bs.
\end{enumerate}
\end{ejemplo}

\section{Ejercicios Propuestos de la Unidad}

\subsection*{Nivel 1: Conceptos y Coeficiente de Correlación}
\begin{enumerate}
    \item Explique detalladamente la diferencia conceptual y práctica entre el Coeficiente de Correlación ($r$) y el Coeficiente de Determinación ($R^2$).
    \item Dados los siguientes coeficientes de correlación muestrales obtenidos en diversos estudios de mercado, clasifique cada uno según su fuerza y sentido:
    \begin{enumerate}
        \item $r = +0.87$
        \item $r = -0.94$
        \item $r = +0.12$
        \item $r = -0.55$
        \item $r = 0.00$
    \end{enumerate}
    \item Para una muestra de $n = 10$ empresas, se calcularon las siguientes sumatorias entre inversión en tecnología ($X$) y rentabilidad operativa ($Y$):
    $\sum x = 50$, $\sum y = 80$, $\sum x^2 = 300$, $\sum y^2 = 720$, $\sum xy = 440$.
    Calcule el coeficiente de correlación de Pearson $r$ e interprete su valor.
\end{enumerate}

\subsection*{Nivel 2: Regresión y Análisis de Costos Empresariales}
\begin{enumerate}[resume]
    \item Una empresa de transporte y logística desea modelar el gasto semanal en combustible ($Y$, en miles de Bs) en función del kilometraje recorrido ($X$, en miles de km). Se registraron los datos de 8 camiones:
    \begin{center}
    \begin{tabular}{l|cccccccc}
    \toprule
    Camión & 1 & 2 & 3 & 4 & 5 & 6 & 7 & 8 \\
    \midrule
    Kilometraje ($X$) & 1.2 & 1.8 & 2.0 & 2.5 & 3.0 & 3.5 & 4.0 & 4.8 \\
    Gasto Combustible ($Y$) & 2.8 & 3.9 & 4.2 & 5.1 & 6.0 & 6.8 & 7.9 & 9.2 \\
    \bottomrule
    \end{tabular}
    \end{center}
    \begin{enumerate}
        \item Determine la ecuación de regresión lineal simple mediante mínimos cuadrados.
        \item Interprete el costo fijo base y el costo marginal por kilómetro recorrido.
        \item Calcule el coeficiente de determinación $R^2$ y comente la calidad del ajuste.
        \item Estime el gasto en combustible si un camión recorre $3.2$ miles de kilómetros en una semana.
    \end{enumerate}

    \item En un estudio sobre gestión de cobranzas en una casa comercial, se recolectó el número de llamadas de recordatorio realizadas a clientes morosos ($X$) y los días de retraso en la cancelación de la cuota ($Y$):
    $(1, 35), (2, 28), (3, 24), (4, 18), (5, 12), (6, 10)$.
    \begin{enumerate}
        \item Calcule el coeficiente de correlación $r$.
        \item Ajuste la recta de regresión $\hat{Y} = b_0 + b_1 X$.
        \item Pronostique los días de retraso esperados si se efectúan 4 llamadas de cobranza.
    \end{enumerate}
\end{enumerate}

\subsection*{Nivel 3: Casos Integrales de Pronóstico Financiero}
\begin{enumerate}[resume]
    \item La gerencia financiera de una cadena hotelera en Santa Cruz analizó la relación entre el porcentaje de ocupación mensual ($X$, en \%) y la utilidad operativa neta ($Y$, en decenas de miles de Bs) durante los últimos 10 meses:
    $\sum x = 650$, $\sum y = 480$, $\sum x^2 = 44500$, $\sum y^2 = 24800$, $\sum xy = 32800$.
    \begin{enumerate}
        \item Determine la ecuación de regresión lineal.
        \item Si en temporada alta se proyecta una ocupación del 85\%, estime la utilidad operativa neta esperada.
        \item Calcule e interprete el coeficiente de determinación $R^2$.
    \end{enumerate}

    \item Un auditor fiscal examina la relación entre los ingresos brutos declarados ($X$, en millones de Bs) y el monto total de crédito fiscal IVA computado ($Y$, en cientos de miles de Bs) para un sector comercial:
    $\bar{x} = 12.0$, $\bar{y} = 8.5$, $s_x = 3.2$, $s_y = 2.4$, $r = 0.85$, $n = 30$.
    \begin{enumerate}
        \item Calcule los coeficientes de la recta de regresión lineal $\hat{Y} = b_0 + b_1 X$ utilizando las propiedades de las medias y desviaciones estándar:
        \[
        b_1 = r \cdot \frac{s_y}{s_x}, \quad b_0 = \bar{y} - b_1 \bar{x}
        \]
        \item Si una empresa del sector declara ingresos brutos por $15.0$ millones de Bs, ¿cuál es el crédito fiscal estimado que el auditor esperaría encontrar en sus libros?
        \item Si la empresa declara un crédito fiscal de $14.0$ cientos de miles de Bs para ese nivel de ingresos, ¿qué sospecha justificada podría formular el auditor para iniciar una fiscalización a fondo?
    \end{enumerate}
\end{enumerate}


% --- PARTE FINAL (BACKMATTER) ---
\backmatter

% Apéndice de Tablas
\chapter*{Apéndice: Tablas Estadísticas de Referencia}
\addcontentsline{toc}{chapter}{Apéndice: Tablas Estadísticas de Referencia}
\label{apendice:tablas}

\section*{Tabla A.1: Áreas bajo la Curva Normal Estándar $P(Z \le z)$}

\begin{center}
\small
\begin{tabular}{c|cccccccccc}
\toprule
$z$ & \textbf{0.00} & \textbf{0.01} & \textbf{0.02} & \textbf{0.03} & \textbf{0.04} & \textbf{0.05} & \textbf{0.06} & \textbf{0.07} & \textbf{0.08} & \textbf{0.09} \\
\midrule
\textbf{0.0} & 0.5000 & 0.5040 & 0.5080 & 0.5120 & 0.5160 & 0.5199 & 0.5239 & 0.5279 & 0.5319 & 0.5359 \\
\textbf{0.1} & 0.5398 & 0.5438 & 0.5478 & 0.5517 & 0.5557 & 0.5596 & 0.5636 & 0.5675 & 0.5714 & 0.5753 \\
\textbf{0.2} & 0.5793 & 0.5832 & 0.5871 & 0.5910 & 0.5948 & 0.5987 & 0.6026 & 0.6064 & 0.6103 & 0.6141 \\
\textbf{0.3} & 0.6179 & 0.6217 & 0.6255 & 0.6293 & 0.6331 & 0.6368 & 0.6406 & 0.6443 & 0.6480 & 0.6517 \\
\textbf{0.4} & 0.6554 & 0.6591 & 0.6628 & 0.6664 & 0.6700 & 0.6736 & 0.6772 & 0.6808 & 0.6844 & 0.6879 \\
\textbf{0.5} & 0.6915 & 0.6950 & 0.6985 & 0.7019 & 0.7054 & 0.7088 & 0.7123 & 0.7157 & 0.7190 & 0.7224 \\
\textbf{0.6} & 0.7257 & 0.7291 & 0.7324 & 0.7357 & 0.7389 & 0.7422 & 0.7454 & 0.7486 & 0.7517 & 0.7549 \\
\textbf{0.7} & 0.7580 & 0.7611 & 0.7642 & 0.7673 & 0.7704 & 0.7734 & 0.7764 & 0.7794 & 0.7823 & 0.7852 \\
\textbf{0.8} & 0.7881 & 0.7910 & 0.7939 & 0.7967 & 0.7995 & 0.8023 & 0.8051 & 0.8078 & 0.8106 & 0.8133 \\
\textbf{0.9} & 0.8159 & 0.8186 & 0.8212 & 0.8238 & 0.8264 & 0.8289 & 0.8315 & 0.8340 & 0.8365 & 0.8389 \\
\textbf{1.0} & 0.8413 & 0.8438 & 0.8461 & 0.8485 & 0.8508 & 0.8531 & 0.8554 & 0.8577 & 0.8599 & 0.8621 \\
\textbf{1.1} & 0.8643 & 0.8665 & 0.8686 & 0.8708 & 0.8729 & 0.8749 & 0.8770 & 0.8790 & 0.8810 & 0.8830 \\
\textbf{1.2} & 0.8849 & 0.8869 & 0.8888 & 0.8907 & 0.8925 & 0.8944 & 0.8962 & 0.8980 & 0.8997 & 0.9015 \\
\textbf{1.3} & 0.9032 & 0.9049 & 0.9066 & 0.9082 & 0.9099 & 0.9115 & 0.9131 & 0.9147 & 0.9162 & 0.9177 \\
\textbf{1.4} & 0.9192 & 0.9207 & 0.9222 & 0.9236 & 0.9251 & 0.9265 & 0.9279 & 0.9292 & 0.9306 & 0.9319 \\
\textbf{1.5} & 0.9332 & 0.9345 & 0.9357 & 0.9370 & 0.9382 & 0.9394 & 0.9406 & 0.9418 & 0.9429 & 0.9441 \\
\textbf{1.6} & 0.9452 & 0.9463 & 0.9474 & 0.9484 & 0.9495 & 0.9505 & 0.9515 & 0.9525 & 0.9535 & 0.9545 \\
\textbf{1.7} & 0.9554 & 0.9564 & 0.9573 & 0.9582 & 0.9591 & 0.9599 & 0.9608 & 0.9616 & 0.9625 & 0.9633 \\
\textbf{1.8} & 0.9641 & 0.9649 & 0.9656 & 0.9664 & 0.9671 & 0.9678 & 0.9686 & 0.9693 & 0.9699 & 0.9706 \\
\textbf{1.9} & 0.9713 & 0.9719 & 0.9726 & 0.9732 & 0.9738 & 0.9744 & 0.9750 & 0.9756 & 0.9761 & 0.9767 \\
\textbf{2.0} & 0.9772 & 0.9778 & 0.9783 & 0.9788 & 0.9793 & 0.9798 & 0.9803 & 0.9808 & 0.9812 & 0.9817 \\
\textbf{2.5} & 0.9938 & 0.9940 & 0.9941 & 0.9943 & 0.9945 & 0.9946 & 0.9948 & 0.9949 & 0.9951 & 0.9952 \\
\textbf{3.0} & 0.9987 & 0.9987 & 0.9987 & 0.9988 & 0.9988 & 0.9989 & 0.9989 & 0.9989 & 0.9990 & 0.9990 \\
\bottomrule
\end{tabular}
\end{center}

\vspace{0.8cm}

\section*{Tabla A.2: Valores Críticos de la Distribución $t$ de Student}

\begin{center}
\small
\begin{tabular}{c|ccccc}
\toprule
\textbf{Grados de Libertad} & \textbf{Confianza 90\%} & \textbf{Confianza 95\%} & \textbf{Confianza 98\%} & \textbf{Confianza 99\%} \\
($\nu = n-1$) & ($\alpha/2 = 0.05$) & ($\alpha/2 = 0.025$) & ($\alpha/2 = 0.01$) & ($\alpha/2 = 0.005$) \\
\midrule
\textbf{1}  & 6.314 & 12.706 & 31.821 & 63.657 \\
\textbf{2}  & 2.920 & 4.303  & 6.965  & 9.925  \\
\textbf{3}  & 2.353 & 3.182  & 4.541  & 5.841  \\
\textbf{4}  & 2.132 & 2.776  & 3.747  & 4.604  \\
\textbf{5}  & 2.015 & 2.571  & 3.365  & 4.032  \\
\textbf{6}  & 1.943 & 2.447  & 3.143  & 3.707  \\
\textbf{7}  & 1.895 & 2.365  & 2.998  & 3.499  \\
\textbf{8}  & 1.860 & 2.306  & 2.896  & 3.355  \\
\textbf{9}  & 1.833 & 2.262  & 2.821  & 3.250  \\
\textbf{10} & 1.812 & 2.228  & 2.764  & 3.169  \\
\textbf{12} & 1.782 & 2.179  & 2.681  & 3.055  \\
\textbf{15} & 1.753 & 2.131  & 2.602  & 2.947  \\
\textbf{20} & 1.725 & 2.086  & 2.528  & 2.845  \\
\textbf{25} & 1.708 & 2.060  & 2.485  & 2.787  \\
\textbf{30} & 1.697 & 2.042  & 2.457  & 2.750  \\
$\boldsymbol{\infty}$ ($Z$) & \textbf{1.645} & \textbf{1.960} & \textbf{2.326} & \textbf{2.576} \\
\bottomrule
\end{tabular}
\end{center}


% Bibliografía General
\chapter*{Bibliografía de Referencia}
\addcontentsline{toc}{chapter}{Bibliografía de Referencia}

\begin{enumerate}
    \item \textbf{Anderson, D. R., Sweeney, D. J., Williams, T. A., Camm, J. D., \& Cochran, J. J.} (2019). \textit{Estadística para administración y economía} (13a ed.). Cengage Learning Editores.
    \item \textbf{Lind, D. A., Marchal, W. G., \& Wathen, S. A.} (2019). \textit{Estadística aplicada a los negocios y la economía} (17a ed.). McGraw-Hill Interamericana.
    \item \textbf{Kazmier, L. J.} (2006). \textit{Estadística aplicada a la administración y economía} (4a ed.). McGraw-Hill (Serie Schaum).
    \item \textbf{Martínez Bencardino, C.} (2012). \textit{Estadística y muestreo} (13a ed.). Ecoe Ediciones.
    \item \textbf{Moya Calderón, R.} (2014). \textit{Probabilidad e inferencia estadística}. Editorial San Marcos.
    \item \textbf{Devore, J. L.} (2016). \textit{Probabilidad y estadística para ingeniería y ciencias} (9a ed.). Cengage Learning.
    \item \textbf{Walpole, R. E., Myers, R. H., Myers, S. L., \& Ye, K.} (2012). \textit{Probabilidad y estadística para ingeniería y ciencias} (9a ed.). Pearson Educación.
    \item \textbf{Zylberberg, A. D.} (2006). \textit{Probabilidad y estadística}. Editorial Nueva Librería.
\end{enumerate}

\end{document}
